% Options for packages loaded elsewhere
\PassOptionsToPackage{unicode}{hyperref}
\PassOptionsToPackage{hyphens}{url}
%
\documentclass[
  ignorenonframetext,
  aspectratio=169,
]{beamer}
\usepackage{pgfpages}
\setbeamertemplate{caption}[numbered]
\setbeamertemplate{caption label separator}{: }
\setbeamercolor{caption name}{fg=normal text.fg}
\beamertemplatenavigationsymbolsempty
% Prevent slide breaks in the middle of a paragraph
\widowpenalties 1 10000
\raggedbottom
\setbeamertemplate{part page}{
  \centering
  \begin{beamercolorbox}[sep=16pt,center]{part title}
    \usebeamerfont{part title}\insertpart\par
  \end{beamercolorbox}
}
\setbeamertemplate{section page}{
  \centering
  \begin{beamercolorbox}[sep=12pt,center]{part title}
    \usebeamerfont{section title}\insertsection\par
  \end{beamercolorbox}
}
\setbeamertemplate{subsection page}{
  \centering
  \begin{beamercolorbox}[sep=8pt,center]{part title}
    \usebeamerfont{subsection title}\insertsubsection\par
  \end{beamercolorbox}
}
\AtBeginPart{
  \frame{\partpage}
}
\AtBeginSection{
  \ifbibliography
  \else
    \frame{\sectionpage}
  \fi
}
\AtBeginSubsection{
  \frame{\subsectionpage}
}

\usepackage{amsmath,amssymb}
\usepackage{iftex}
\ifPDFTeX
  \usepackage[T1]{fontenc}
  \usepackage[utf8]{inputenc}
  \usepackage{textcomp} % provide euro and other symbols
\else % if luatex or xetex
  \usepackage{unicode-math}
  \defaultfontfeatures{Scale=MatchLowercase}
  \defaultfontfeatures[\rmfamily]{Ligatures=TeX,Scale=1}
\fi
\usepackage{lmodern}
\usetheme[]{Madrid}
\usecolortheme{default}
\ifPDFTeX\else  
    % xetex/luatex font selection
\fi
% Use upquote if available, for straight quotes in verbatim environments
\IfFileExists{upquote.sty}{\usepackage{upquote}}{}
\IfFileExists{microtype.sty}{% use microtype if available
  \usepackage[]{microtype}
  \UseMicrotypeSet[protrusion]{basicmath} % disable protrusion for tt fonts
}{}
\makeatletter
\@ifundefined{KOMAClassName}{% if non-KOMA class
  \IfFileExists{parskip.sty}{%
    \usepackage{parskip}
  }{% else
    \setlength{\parindent}{0pt}
    \setlength{\parskip}{6pt plus 2pt minus 1pt}}
}{% if KOMA class
  \KOMAoptions{parskip=half}}
\makeatother
\usepackage{xcolor}
\newif\ifbibliography
\setlength{\emergencystretch}{3em} % prevent overfull lines
\setcounter{secnumdepth}{5}

\usepackage{color}
\usepackage{fancyvrb}
\newcommand{\VerbBar}{|}
\newcommand{\VERB}{\Verb[commandchars=\\\{\}]}
\DefineVerbatimEnvironment{Highlighting}{Verbatim}{commandchars=\\\{\}}
% Add ',fontsize=\small' for more characters per line
\usepackage{framed}
\definecolor{shadecolor}{RGB}{241,243,245}
\newenvironment{Shaded}{\begin{snugshade}}{\end{snugshade}}
\newcommand{\AlertTok}[1]{\textcolor[rgb]{0.68,0.00,0.00}{#1}}
\newcommand{\AnnotationTok}[1]{\textcolor[rgb]{0.37,0.37,0.37}{#1}}
\newcommand{\AttributeTok}[1]{\textcolor[rgb]{0.40,0.45,0.13}{#1}}
\newcommand{\BaseNTok}[1]{\textcolor[rgb]{0.68,0.00,0.00}{#1}}
\newcommand{\BuiltInTok}[1]{\textcolor[rgb]{0.00,0.23,0.31}{#1}}
\newcommand{\CharTok}[1]{\textcolor[rgb]{0.13,0.47,0.30}{#1}}
\newcommand{\CommentTok}[1]{\textcolor[rgb]{0.37,0.37,0.37}{#1}}
\newcommand{\CommentVarTok}[1]{\textcolor[rgb]{0.37,0.37,0.37}{\textit{#1}}}
\newcommand{\ConstantTok}[1]{\textcolor[rgb]{0.56,0.35,0.01}{#1}}
\newcommand{\ControlFlowTok}[1]{\textcolor[rgb]{0.00,0.23,0.31}{#1}}
\newcommand{\DataTypeTok}[1]{\textcolor[rgb]{0.68,0.00,0.00}{#1}}
\newcommand{\DecValTok}[1]{\textcolor[rgb]{0.68,0.00,0.00}{#1}}
\newcommand{\DocumentationTok}[1]{\textcolor[rgb]{0.37,0.37,0.37}{\textit{#1}}}
\newcommand{\ErrorTok}[1]{\textcolor[rgb]{0.68,0.00,0.00}{#1}}
\newcommand{\ExtensionTok}[1]{\textcolor[rgb]{0.00,0.23,0.31}{#1}}
\newcommand{\FloatTok}[1]{\textcolor[rgb]{0.68,0.00,0.00}{#1}}
\newcommand{\FunctionTok}[1]{\textcolor[rgb]{0.28,0.35,0.67}{#1}}
\newcommand{\ImportTok}[1]{\textcolor[rgb]{0.00,0.46,0.62}{#1}}
\newcommand{\InformationTok}[1]{\textcolor[rgb]{0.37,0.37,0.37}{#1}}
\newcommand{\KeywordTok}[1]{\textcolor[rgb]{0.00,0.23,0.31}{#1}}
\newcommand{\NormalTok}[1]{\textcolor[rgb]{0.00,0.23,0.31}{#1}}
\newcommand{\OperatorTok}[1]{\textcolor[rgb]{0.37,0.37,0.37}{#1}}
\newcommand{\OtherTok}[1]{\textcolor[rgb]{0.00,0.23,0.31}{#1}}
\newcommand{\PreprocessorTok}[1]{\textcolor[rgb]{0.68,0.00,0.00}{#1}}
\newcommand{\RegionMarkerTok}[1]{\textcolor[rgb]{0.00,0.23,0.31}{#1}}
\newcommand{\SpecialCharTok}[1]{\textcolor[rgb]{0.37,0.37,0.37}{#1}}
\newcommand{\SpecialStringTok}[1]{\textcolor[rgb]{0.13,0.47,0.30}{#1}}
\newcommand{\StringTok}[1]{\textcolor[rgb]{0.13,0.47,0.30}{#1}}
\newcommand{\VariableTok}[1]{\textcolor[rgb]{0.07,0.07,0.07}{#1}}
\newcommand{\VerbatimStringTok}[1]{\textcolor[rgb]{0.13,0.47,0.30}{#1}}
\newcommand{\WarningTok}[1]{\textcolor[rgb]{0.37,0.37,0.37}{\textit{#1}}}

\providecommand{\tightlist}{%
  \setlength{\itemsep}{0pt}\setlength{\parskip}{0pt}}\usepackage{longtable,booktabs,array}
\usepackage{calc} % for calculating minipage widths
\usepackage{caption}
% Make caption package work with longtable
\makeatletter
\def\fnum@table{\tablename~\thetable}
\makeatother
\usepackage{graphicx}
\makeatletter
\def\maxwidth{\ifdim\Gin@nat@width>\linewidth\linewidth\else\Gin@nat@width\fi}
\def\maxheight{\ifdim\Gin@nat@height>\textheight\textheight\else\Gin@nat@height\fi}
\makeatother
% Scale images if necessary, so that they will not overflow the page
% margins by default, and it is still possible to overwrite the defaults
% using explicit options in \includegraphics[width, height, ...]{}
\setkeys{Gin}{width=\maxwidth,height=\maxheight,keepaspectratio}
% Set default figure placement to htbp
\makeatletter
\def\fps@figure{htbp}
\makeatother

\usepackage{etoolbox}
% Smaller font for code and output so they fit on slides
\AtBeginEnvironment{Shaded}{\scriptsize}
\makeatletter
\def\verbatim@font{\scriptsize\ttfamily}
\makeatother
% Outline slide at the start of each section
\AtBeginSection[]{\begin{frame}{Outline}\tableofcontents[currentsection]\end{frame}}
\makeatletter
\makeatother
\makeatletter
\makeatother
\makeatletter
\@ifpackageloaded{caption}{}{\usepackage{caption}}
\AtBeginDocument{%
\ifdefined\contentsname
  \renewcommand*\contentsname{Table of contents}
\else
  \newcommand\contentsname{Table of contents}
\fi
\ifdefined\listfigurename
  \renewcommand*\listfigurename{List of Figures}
\else
  \newcommand\listfigurename{List of Figures}
\fi
\ifdefined\listtablename
  \renewcommand*\listtablename{List of Tables}
\else
  \newcommand\listtablename{List of Tables}
\fi
\ifdefined\figurename
  \renewcommand*\figurename{Figure}
\else
  \newcommand\figurename{Figure}
\fi
\ifdefined\tablename
  \renewcommand*\tablename{Table}
\else
  \newcommand\tablename{Table}
\fi
}
\@ifpackageloaded{float}{}{\usepackage{float}}
\floatstyle{ruled}
\@ifundefined{c@chapter}{\newfloat{codelisting}{h}{lop}}{\newfloat{codelisting}{h}{lop}[chapter]}
\floatname{codelisting}{Listing}
\newcommand*\listoflistings{\listof{codelisting}{List of Listings}}
\makeatother
\makeatletter
\@ifpackageloaded{caption}{}{\usepackage{caption}}
\@ifpackageloaded{subcaption}{}{\usepackage{subcaption}}
\makeatother
\makeatletter
\@ifpackageloaded{tcolorbox}{}{\usepackage[skins,breakable]{tcolorbox}}
\makeatother
\makeatletter
\@ifundefined{shadecolor}{\definecolor{shadecolor}{rgb}{.97, .97, .97}}
\makeatother
\makeatletter
\makeatother
\makeatletter
\makeatother
\ifLuaTeX
  \usepackage{selnolig}  % disable illegal ligatures
\fi
\IfFileExists{bookmark.sty}{\usepackage{bookmark}}{\usepackage{hyperref}}
\IfFileExists{xurl.sty}{\usepackage{xurl}}{} % add URL line breaks if available
\urlstyle{same} % disable monospaced font for URLs
\hypersetup{
  pdftitle={Session 2: Data Types and Structures},
  pdfauthor={Taikgun Song},
  hidelinks,
  pdfcreator={LaTeX via pandoc}}

\title{Session 2: Data Types and Structures}
\subtitle{R Workshop for DBA Students}
\author{Taikgun Song}
\date{2026-11-06}

\begin{document}
\frame{\titlepage}
\ifdefined\Shaded\renewenvironment{Shaded}{\begin{tcolorbox}[breakable, sharp corners, enhanced, frame hidden, borderline west={3pt}{0pt}{shadecolor}, interior hidden, boxrule=0pt]}{\end{tcolorbox}}\fi

\renewcommand*\contentsname{Outline}
\begin{frame}[allowframebreaks]
  \frametitle{Outline}
  \tableofcontents[hideallsubsections]
\end{frame}
\begin{frame}
\end{frame}

\begin{frame}{Session Objectives}
\protect\hypertarget{session-objectives}{}
By the end of this session, you will be able to:

\begin{enumerate}
\tightlist
\item
  Explain why the way a variable is \emph{stored} determines how R, and
  your statistical models, treat it
\item
  Identify R's basic data types and the special types built on them
  (factor, Date, date-time)
\item
  Create and modify \textbf{factors}, including applying a data
  dictionary and setting a reference level
\item
  Read regression output for a categorical predictor under reference
  coding and one-hot coding
\item
  Create, subset, and modify R's data structures: vector, matrix, array,
  list, and data frame
\end{enumerate}
\end{frame}

\begin{frame}[fragile]{Before You Start}
\protect\hypertarget{before-you-start}{}
\begin{itemize}
\tightlist
\item
  Open your workshop RStudio Project (Session 1)
\item
  Load the packages for today:
\end{itemize}

\begin{Shaded}
\begin{Highlighting}[]
\FunctionTok{library}\NormalTok{(tidyverse)   }\CommentTok{\# includes dplyr, ggplot2, forcats, lubridate}
\FunctionTok{library}\NormalTok{(carData)     }\CommentTok{\# installed together with the car package}
\end{Highlighting}
\end{Shaded}

\begin{itemize}
\tightlist
\item
  If you get an error, install them first:
\end{itemize}

\begin{Shaded}
\begin{Highlighting}[]
\FunctionTok{install.packages}\NormalTok{(}\FunctionTok{c}\NormalTok{(}\StringTok{"tidyverse"}\NormalTok{, }\StringTok{"car"}\NormalTok{))}
\end{Highlighting}
\end{Shaded}

\begin{itemize}
\tightlist
\item
  This session uses \texttt{\textless{}-} to store values
  (\texttt{x\ \textless{}-\ 5}) and comparisons such as \texttt{==} and
  \texttt{\textgreater{}}. \textbf{Session 3} covers operators in detail
\end{itemize}
\end{frame}

\hypertarget{sec-why}{%
\section{Why Data Types Matter}\label{sec-why}}

\begin{frame}{Excel vs.~R}
\protect\hypertarget{excel-vs.-r}{}
\begin{itemize}
\tightlist
\item
  In Excel, a cell can hold anything, and Excel guesses what you meant
\item
  In R, every value has a \textbf{type}, and the type decides what you
  can do with it
\item
  Getting the type right is not a technicality: \textbf{it changes the
  answers your models give you}
\end{itemize}
\end{frame}

\begin{frame}{Numbers That Are Not Really Numbers}
\protect\hypertarget{sec-coded}{}
\begin{columns}[T]
\begin{column}{0.55\textwidth}
\begin{itemize}
\tightlist
\item
  Census, government surveys, and corporate databases often store
  categories as \textbf{numeric codes} to save space
\item
  A \textbf{data dictionary} (codebook) says what each code means
\item
  The dictionary is for humans. R only sees the numbers
\end{itemize}
\end{column}

\begin{column}{0.45\textwidth}
\begin{longtable}[]{@{}ll@{}}
\toprule\noalign{}
Code & Department \\
\midrule\noalign{}
\endhead
1 & Sales \\
2 & HR \\
3 & IT \\
\bottomrule\noalign{}
\end{longtable}
\end{column}
\end{columns}
\end{frame}

\begin{frame}[fragile]{What R Sees}
\protect\hypertarget{what-r-sees}{}
R reads \texttt{department} as an \textbf{integer} and happily does
arithmetic with it:

\begin{Shaded}
\begin{Highlighting}[]
\NormalTok{department }\OtherTok{\textless{}{-}} \FunctionTok{c}\NormalTok{(}\DecValTok{1}\NormalTok{, }\DecValTok{3}\NormalTok{, }\DecValTok{2}\NormalTok{, }\DecValTok{1}\NormalTok{, }\DecValTok{1}\NormalTok{, }\DecValTok{3}\NormalTok{, }\DecValTok{2}\NormalTok{, }\DecValTok{3}\NormalTok{)}

\FunctionTok{typeof}\NormalTok{(department)}
\end{Highlighting}
\end{Shaded}

\begingroup\tiny\color{darkgray}

\begin{verbatim}
[1] "double"
\end{verbatim}

\endgroup

\begin{Shaded}
\begin{Highlighting}[]
\FunctionTok{mean}\NormalTok{(department)                }\CommentTok{\# the "average department"?}
\end{Highlighting}
\end{Shaded}

\begingroup\tiny\color{darkgray}

\begin{verbatim}
[1] 2
\end{verbatim}

\endgroup

\begin{Shaded}
\begin{Highlighting}[]
\NormalTok{department[}\DecValTok{1}\NormalTok{] }\SpecialCharTok{\textless{}}\NormalTok{ department[}\DecValTok{2}\NormalTok{]   }\CommentTok{\# is Sales "less than" IT?}
\end{Highlighting}
\end{Shaded}

\begingroup\tiny\color{darkgray}

\begin{verbatim}
[1] TRUE
\end{verbatim}

\endgroup
\end{frame}

\begin{frame}[fragile]{Why It Matters for Regression}
\protect\hypertarget{why-it-matters-for-regression}{}
If \texttt{department} stays numeric, a regression estimates \textbf{one
slope} for it:

\begin{itemize}
\tightlist
\item
  It assumes HR is ``one unit more'' than Sales, and IT ``two units
  more''
\item
  But 1, 2, 3 were \textbf{arbitrary labels}. Code IT as 1 and the model
  gives a different answer
\item
  Underneath, \texttt{department} is a category, not a quantity
\end{itemize}

\begin{block}{The fix}
\protect\hypertarget{the-fix}{}
Convert it to a \textbf{factor} and apply the data dictionary (next
slide). More on factors in Section~\ref{sec-factors} and on regression
in Section~\ref{sec-lm}.
\end{block}
\end{frame}

\begin{frame}[fragile]{Apply the Data Dictionary}
\protect\hypertarget{apply-the-data-dictionary}{}
\begin{Shaded}
\begin{Highlighting}[]
\NormalTok{department }\OtherTok{\textless{}{-}} \FunctionTok{factor}\NormalTok{(department,}
                     \AttributeTok{levels =} \FunctionTok{c}\NormalTok{(}\DecValTok{1}\NormalTok{, }\DecValTok{2}\NormalTok{, }\DecValTok{3}\NormalTok{),}
                     \AttributeTok{labels =} \FunctionTok{c}\NormalTok{(}\StringTok{"Sales"}\NormalTok{, }\StringTok{"HR"}\NormalTok{, }\StringTok{"IT"}\NormalTok{))}
\NormalTok{department}
\end{Highlighting}
\end{Shaded}

\begingroup\tiny\color{darkgray}

\begin{verbatim}
[1] Sales IT    HR    Sales Sales IT    HR    IT   
Levels: Sales HR IT
\end{verbatim}

\endgroup

\begin{Shaded}
\begin{Highlighting}[]
\FunctionTok{table}\NormalTok{(department)}
\end{Highlighting}
\end{Shaded}

\begingroup\tiny\color{darkgray}

\begin{verbatim}
department
Sales    HR    IT 
    3     2     3 
\end{verbatim}

\endgroup
\end{frame}

\begin{frame}{Likert Items: A Gray Area}
\protect\hypertarget{sec-likert}{}
\begin{columns}[T]
\begin{column}{0.5\textwidth}
\emph{``I am satisfied with my manager''}

\begin{longtable}[]{@{}ll@{}}
\toprule\noalign{}
Code & Response \\
\midrule\noalign{}
\endhead
1 & Strongly disagree \\
2 & Disagree \\
3 & Neutral \\
4 & Agree \\
5 & Strongly agree \\
\bottomrule\noalign{}
\end{longtable}
\end{column}

\begin{column}{0.5\textwidth}
\begin{itemize}
\tightlist
\item
  Unlike department codes, these numbers \textbf{do} have an order
\item
  But is the distance from 1 to 2 the same as from 4 to 5?
\item
  \textbf{Nobody knows.} So there is more than one reasonable way to
  store the item
\end{itemize}
\end{column}
\end{columns}
\end{frame}

\begin{frame}{Three Ways to Store a Likert Item}
\protect\hypertarget{three-ways-to-store-a-likert-item}{}
\begin{longtable}[]{@{}
  >{\raggedright\arraybackslash}p{(\columnwidth - 4\tabcolsep) * \real{0.3824}}
  >{\raggedright\arraybackslash}p{(\columnwidth - 4\tabcolsep) * \real{0.2647}}
  >{\raggedright\arraybackslash}p{(\columnwidth - 4\tabcolsep) * \real{0.3529}}@{}}
\toprule\noalign{}
\begin{minipage}[b]{\linewidth}\raggedright
Store it as
\end{minipage} & \begin{minipage}[b]{\linewidth}\raggedright
Assumes
\end{minipage} & \begin{minipage}[b]{\linewidth}\raggedright
Common use
\end{minipage} \\
\midrule\noalign{}
\endhead
Numeric (1 to 5) & Equal spacing & Averaging items into a scale score \\
Ordered factor & Order, unknown spacing & Sorting, tables, plots,
ordinal models \\
Unordered factor & Separate categories & Comparing specific groups \\
\bottomrule\noalign{}
\end{longtable}

\begin{block}{Key points}
\protect\hypertarget{key-points}{}
\begin{itemize}
\tightlist
\item
  \textbf{Know which choice you made.} R will not warn you either way
\item
  \textbf{Multi-item scales} (e.g.~averaging five items) are usually
  treated as numeric
\item
  \textbf{Ordered factors behave unexpectedly in regression} (see
  Section~\ref{sec-ordered})
\end{itemize}
\end{block}
\end{frame}

\begin{frame}[fragile]{Likert Items in R}
\protect\hypertarget{likert-items-in-r}{}
\begin{Shaded}
\begin{Highlighting}[]
\NormalTok{satisfaction }\OtherTok{\textless{}{-}} \FunctionTok{c}\NormalTok{(}\DecValTok{4}\NormalTok{, }\DecValTok{5}\NormalTok{, }\DecValTok{2}\NormalTok{, }\DecValTok{4}\NormalTok{, }\DecValTok{3}\NormalTok{, }\DecValTok{1}\NormalTok{, }\DecValTok{5}\NormalTok{, }\DecValTok{4}\NormalTok{)}
\FunctionTok{mean}\NormalTok{(satisfaction)          }\CommentTok{\# as numeric: assumes equal spacing}
\end{Highlighting}
\end{Shaded}

\begingroup\tiny\color{darkgray}

\begin{verbatim}
[1] 3.5
\end{verbatim}

\endgroup

\begin{Shaded}
\begin{Highlighting}[]
\NormalTok{satisfaction\_f }\OtherTok{\textless{}{-}} \FunctionTok{factor}\NormalTok{(satisfaction, }\AttributeTok{levels =} \DecValTok{1}\SpecialCharTok{:}\DecValTok{5}\NormalTok{,}
                         \AttributeTok{labels =} \FunctionTok{c}\NormalTok{(}\StringTok{"Strongly disagree"}\NormalTok{, }\StringTok{"Disagree"}\NormalTok{, }\StringTok{"Neutral"}\NormalTok{,}
                                    \StringTok{"Agree"}\NormalTok{, }\StringTok{"Strongly agree"}\NormalTok{),}
                         \AttributeTok{ordered =} \ConstantTok{TRUE}\NormalTok{)}
\NormalTok{satisfaction\_f}
\end{Highlighting}
\end{Shaded}

\begingroup\tiny\color{darkgray}

\begin{verbatim}
[1] Agree             Strongly agree    Disagree         
[4] Agree             Neutral           Strongly disagree
[7] Strongly agree    Agree            
5 Levels: Strongly disagree < Disagree < Neutral < ... < Strongly agree
\end{verbatim}

\endgroup

\begin{Shaded}
\begin{Highlighting}[]
\NormalTok{satisfaction\_f }\SpecialCharTok{\textgreater{}=} \StringTok{"Agree"}   \CommentTok{\# comparisons work because the order is defined}
\end{Highlighting}
\end{Shaded}

\begingroup\tiny\color{darkgray}

\begin{verbatim}
[1]  TRUE  TRUE FALSE  TRUE FALSE FALSE  TRUE  TRUE
\end{verbatim}

\endgroup
\end{frame}

\hypertarget{sec-types}{%
\section{Data Types}\label{sec-types}}

\begin{frame}[fragile]{The Four Basic Types}
\protect\hypertarget{sec-basic-types}{}
\begin{longtable}[]{@{}
  >{\raggedright\arraybackslash}p{(\columnwidth - 4\tabcolsep) * \real{0.1818}}
  >{\raggedright\arraybackslash}p{(\columnwidth - 4\tabcolsep) * \real{0.2727}}
  >{\raggedright\arraybackslash}p{(\columnwidth - 4\tabcolsep) * \real{0.5455}}@{}}
\toprule\noalign{}
\begin{minipage}[b]{\linewidth}\raggedright
Type
\end{minipage} & \begin{minipage}[b]{\linewidth}\raggedright
Example
\end{minipage} & \begin{minipage}[b]{\linewidth}\raggedright
Business example
\end{minipage} \\
\midrule\noalign{}
\endhead
Logical & \texttt{TRUE}, \texttt{FALSE} & Did the customer renew? (only
two values) \\
Integer & \texttt{5L} & Number of employees \\
Double (numeric) & \texttt{3.14} & Revenue, profit margin \\
Character & \texttt{"Sales"} & Customer name, survey comment \\
\bottomrule\noalign{}
\end{longtable}

\begin{Shaded}
\begin{Highlighting}[]
\NormalTok{renewed }\OtherTok{\textless{}{-}} \ConstantTok{TRUE}\NormalTok{; employees }\OtherTok{\textless{}{-}}\NormalTok{ 250L; revenue }\OtherTok{\textless{}{-}} \FloatTok{1250000.50}\NormalTok{; dept }\OtherTok{\textless{}{-}} \StringTok{"Sales"}
\FunctionTok{c}\NormalTok{(}\FunctionTok{typeof}\NormalTok{(renewed), }\FunctionTok{typeof}\NormalTok{(employees), }\FunctionTok{typeof}\NormalTok{(revenue), }\FunctionTok{typeof}\NormalTok{(dept))}
\end{Highlighting}
\end{Shaded}

\begingroup\tiny\color{darkgray}

\begin{verbatim}
[1] "logical"   "integer"   "double"    "character"
\end{verbatim}

\endgroup
\end{frame}

\begin{frame}[fragile]{Special Types Built on the Basic Types}
\protect\hypertarget{sec-special-types}{}
Each is a basic type with extra information attached, which R calls a
\textbf{class}:

\begin{longtable}[]{@{}
  >{\raggedright\arraybackslash}p{(\columnwidth - 6\tabcolsep) * \real{0.2692}}
  >{\raggedright\arraybackslash}p{(\columnwidth - 6\tabcolsep) * \real{0.1346}}
  >{\raggedright\arraybackslash}p{(\columnwidth - 6\tabcolsep) * \real{0.4231}}
  >{\raggedright\arraybackslash}p{(\columnwidth - 6\tabcolsep) * \real{0.1731}}@{}}
\toprule\noalign{}
\begin{minipage}[b]{\linewidth}\raggedright
Kind of data
\end{minipage} & \begin{minipage}[b]{\linewidth}\raggedright
Class
\end{minipage} & \begin{minipage}[b]{\linewidth}\raggedright
Stored underneath as
\end{minipage} & \begin{minipage}[b]{\linewidth}\raggedright
Example
\end{minipage} \\
\midrule\noalign{}
\endhead
Categorical & \texttt{factor} & Integer codes + labels & Department \\
Date & \texttt{Date} & Days since 1970-01-01 & Hire date \\
Date-time & \texttt{POSIXct} & Seconds since 1970-01-01 UTC & Order
timestamp \\
\bottomrule\noalign{}
\end{longtable}

\begin{Shaded}
\begin{Highlighting}[]
\FunctionTok{typeof}\NormalTok{(department); }\FunctionTok{class}\NormalTok{(department)   }\CommentTok{\# stored as integers, treated as a factor}
\end{Highlighting}
\end{Shaded}

\begingroup\tiny\color{darkgray}

\begin{verbatim}
[1] "integer"
\end{verbatim}

\endgroup

\begingroup\tiny\color{darkgray}

\begin{verbatim}
[1] "factor"
\end{verbatim}

\endgroup

\begin{block}{Notice}
\protect\hypertarget{notice}{}
A factor is stored exactly like the department codes: integers plus a
dictionary of labels. The difference is that R now \emph{knows} they are
codes.
\end{block}
\end{frame}

\begin{frame}[fragile]{Dates}
\protect\hypertarget{dates}{}
\begin{Shaded}
\begin{Highlighting}[]
\NormalTok{hire\_date }\OtherTok{\textless{}{-}} \FunctionTok{as.Date}\NormalTok{(}\StringTok{"2021{-}09{-}23"}\NormalTok{)              }\CommentTok{\# ISO format: year{-}month{-}day}
\FunctionTok{as.Date}\NormalTok{(}\StringTok{"09/23/2021"}\NormalTok{, }\AttributeTok{format =} \StringTok{"\%m/\%d/\%Y"}\NormalTok{)      }\CommentTok{\# other formats need a format string}
\end{Highlighting}
\end{Shaded}

\begingroup\tiny\color{darkgray}

\begin{verbatim}
[1] "2021-09-23"
\end{verbatim}

\endgroup

\begin{Shaded}
\begin{Highlighting}[]
\CommentTok{\# lubridate: name the order of the parts}
\FunctionTok{ymd}\NormalTok{(}\StringTok{"2021{-}09{-}23"}\NormalTok{); }\FunctionTok{mdy}\NormalTok{(}\StringTok{"09/23/2021"}\NormalTok{); }\FunctionTok{dmy}\NormalTok{(}\StringTok{"23.09.2021"}\NormalTok{)}
\end{Highlighting}
\end{Shaded}

\begingroup\tiny\color{darkgray}

\begin{verbatim}
[1] "2021-09-23"
\end{verbatim}

\endgroup

\begingroup\tiny\color{darkgray}

\begin{verbatim}
[1] "2021-09-23"
\end{verbatim}

\endgroup

\begingroup\tiny\color{darkgray}

\begin{verbatim}
[1] "2021-09-23"
\end{verbatim}

\endgroup

\begin{Shaded}
\begin{Highlighting}[]
\NormalTok{as\_of }\OtherTok{\textless{}{-}} \FunctionTok{as.Date}\NormalTok{(}\StringTok{"2026{-}11{-}04"}\NormalTok{)}
\NormalTok{as\_of }\SpecialCharTok{{-}}\NormalTok{ hire\_date                                }\CommentTok{\# tenure in days}
\end{Highlighting}
\end{Shaded}

\begingroup\tiny\color{darkgray}

\begin{verbatim}
Time difference of 1868 days
\end{verbatim}

\endgroup

\begin{Shaded}
\begin{Highlighting}[]
\FunctionTok{year}\NormalTok{(hire\_date); }\FunctionTok{month}\NormalTok{(hire\_date, }\AttributeTok{label =} \ConstantTok{TRUE}\NormalTok{)}
\end{Highlighting}
\end{Shaded}

\begingroup\tiny\color{darkgray}

\begin{verbatim}
[1] 2021
\end{verbatim}

\endgroup

\begingroup\tiny\color{darkgray}

\begin{verbatim}
[1] 9
Levels: 1 < 2 < 3 < 4 < 5 < 6 < 7 < 8 < 9 < 10 < 11 < 12
\end{verbatim}

\endgroup
\end{frame}

\begin{frame}[fragile]{Date-Times (\texttt{POSIXct})}
\protect\hypertarget{date-times-posixct}{}
A date-time stores a moment down to the second, so it also needs a
\textbf{time zone}.

\begin{Shaded}
\begin{Highlighting}[]
\NormalTok{order\_time }\OtherTok{\textless{}{-}} \FunctionTok{ymd\_hms}\NormalTok{(}\StringTok{"2026{-}11{-}04 12:38:00"}\NormalTok{, }\AttributeTok{tz =} \StringTok{"America/New\_York"}\NormalTok{)}
\FunctionTok{class}\NormalTok{(order\_time); }\FunctionTok{typeof}\NormalTok{(order\_time)}
\end{Highlighting}
\end{Shaded}

\begingroup\tiny\color{darkgray}

\begin{verbatim}
[1] "POSIXct" "POSIXt" 
\end{verbatim}

\endgroup

\begingroup\tiny\color{darkgray}

\begin{verbatim}
[1] "double"
\end{verbatim}

\endgroup

\begin{Shaded}
\begin{Highlighting}[]
\FunctionTok{hour}\NormalTok{(order\_time); }\FunctionTok{wday}\NormalTok{(order\_time, }\AttributeTok{label =} \ConstantTok{TRUE}\NormalTok{)}
\end{Highlighting}
\end{Shaded}

\begingroup\tiny\color{darkgray}

\begin{verbatim}
[1] 12
\end{verbatim}

\endgroup

\begingroup\tiny\color{darkgray}

\begin{verbatim}
[1] 수
Levels: 일 < 월 < 화 < 수 < 목 < 금 < 토
\end{verbatim}

\endgroup

\begin{Shaded}
\begin{Highlighting}[]
\NormalTok{ship\_time }\OtherTok{\textless{}{-}} \FunctionTok{ymd\_hms}\NormalTok{(}\StringTok{"2026{-}11{-}05 09:15:00"}\NormalTok{, }\AttributeTok{tz =} \StringTok{"America/New\_York"}\NormalTok{)}
\FunctionTok{difftime}\NormalTok{(ship\_time, order\_time, }\AttributeTok{units =} \StringTok{"hours"}\NormalTok{)}
\end{Highlighting}
\end{Shaded}

\begingroup\tiny\color{darkgray}

\begin{verbatim}
Time difference of 20.61667 hours
\end{verbatim}

\endgroup

\begin{Shaded}
\begin{Highlighting}[]
\FunctionTok{with\_tz}\NormalTok{(order\_time, }\AttributeTok{tzone =} \StringTok{"Asia/Seoul"}\NormalTok{)        }\CommentTok{\# same moment, another zone}
\end{Highlighting}
\end{Shaded}

\begingroup\tiny\color{darkgray}

\begin{verbatim}
[1] "2026-11-05 02:38:00 KST"
\end{verbatim}

\endgroup
\end{frame}

\begin{frame}[fragile]{Time Zones}
\protect\hypertarget{time-zones}{}
\begin{alertblock}{Watch out}
\protect\hypertarget{watch-out}{}
If you do not set \texttt{tz}, R uses your computer's time zone. The
same code can give different results on different computers.

\begin{itemize}
\tightlist
\item
  Base R equivalent:
  \texttt{as.POSIXct("2026-11-04\ 12:38:00",\ tz\ =\ "America/New\_York")}
\item
  Survey platforms such as Qualtrics often export timestamps in UTC or
  in the survey owner's time zone
\item
  Check the documentation and \textbf{set \texttt{tz} explicitly}
\end{itemize}
\end{alertblock}
\end{frame}

\begin{frame}[fragile]{Checking and Converting Types}
\protect\hypertarget{sec-convert}{}
\begin{itemize}
\tightlist
\item
  \texttt{is.*()} functions \textbf{check} a type; \texttt{as.*()}
  functions \textbf{convert} to a type
\end{itemize}

\begin{Shaded}
\begin{Highlighting}[]
\FunctionTok{is.numeric}\NormalTok{(revenue); }\FunctionTok{is.character}\NormalTok{(revenue)}
\end{Highlighting}
\end{Shaded}

\begingroup\tiny\color{darkgray}

\begin{verbatim}
[1] TRUE
\end{verbatim}

\endgroup

\begingroup\tiny\color{darkgray}

\begin{verbatim}
[1] FALSE
\end{verbatim}

\endgroup

\begin{Shaded}
\begin{Highlighting}[]
\FunctionTok{as.character}\NormalTok{(revenue); }\FunctionTok{as.numeric}\NormalTok{(}\StringTok{"42"}\NormalTok{)}
\end{Highlighting}
\end{Shaded}

\begingroup\tiny\color{darkgray}

\begin{verbatim}
[1] "1250000.5"
\end{verbatim}

\endgroup

\begingroup\tiny\color{darkgray}

\begin{verbatim}
[1] 42
\end{verbatim}

\endgroup

\begin{itemize}
\tightlist
\item
  When a conversion is impossible, R returns \texttt{NA} \textbf{with a
  warning}:
\end{itemize}

\begin{Shaded}
\begin{Highlighting}[]
\FunctionTok{as.numeric}\NormalTok{(}\FunctionTok{c}\NormalTok{(}\StringTok{"100"}\NormalTok{, }\StringTok{"250"}\NormalTok{, }\StringTok{"N/A"}\NormalTok{, }\StringTok{"$1,200"}\NormalTok{))}
\end{Highlighting}
\end{Shaded}

\begingroup\tiny\color{darkgray}

\begin{verbatim}
Warning: NAs introduced by coercion
\end{verbatim}

\endgroup

\begingroup\tiny\color{darkgray}

\begin{verbatim}
[1] 100 250  NA  NA
\end{verbatim}

\endgroup

\begin{exampleblock}{Excel users}
\protect\hypertarget{excel-users}{}
A column that \emph{looks} numeric in Excel often arrives in R as
character because a few cells contain text, commas, or currency symbols.
\end{exampleblock}
\end{frame}

\begin{frame}[fragile]{Missing Values}
\protect\hypertarget{sec-na}{}
\texttt{NA} means ``not available.'' Most calculations involving
\texttt{NA} return \texttt{NA}.

\begin{Shaded}
\begin{Highlighting}[]
\NormalTok{scores }\OtherTok{\textless{}{-}} \FunctionTok{c}\NormalTok{(}\DecValTok{4}\NormalTok{, }\DecValTok{5}\NormalTok{, }\ConstantTok{NA}\NormalTok{, }\DecValTok{3}\NormalTok{, }\DecValTok{5}\NormalTok{)}

\FunctionTok{mean}\NormalTok{(scores)                 }\CommentTok{\# NA: R refuses to guess}
\end{Highlighting}
\end{Shaded}

\begingroup\tiny\color{darkgray}

\begin{verbatim}
[1] NA
\end{verbatim}

\endgroup

\begin{Shaded}
\begin{Highlighting}[]
\FunctionTok{mean}\NormalTok{(scores, }\AttributeTok{na.rm =} \ConstantTok{TRUE}\NormalTok{)   }\CommentTok{\# ignore missing values}
\end{Highlighting}
\end{Shaded}

\begingroup\tiny\color{darkgray}

\begin{verbatim}
[1] 4.25
\end{verbatim}

\endgroup

\begin{Shaded}
\begin{Highlighting}[]
\FunctionTok{is.na}\NormalTok{(scores)                }\CommentTok{\# which values are missing?}
\end{Highlighting}
\end{Shaded}

\begingroup\tiny\color{darkgray}

\begin{verbatim}
[1] FALSE FALSE  TRUE FALSE FALSE
\end{verbatim}

\endgroup

\begin{Shaded}
\begin{Highlighting}[]
\FunctionTok{sum}\NormalTok{(}\FunctionTok{is.na}\NormalTok{(scores))           }\CommentTok{\# how many?}
\end{Highlighting}
\end{Shaded}

\begingroup\tiny\color{darkgray}

\begin{verbatim}
[1] 1
\end{verbatim}

\endgroup

\begin{itemize}
\tightlist
\item
  \texttt{scores\ ==\ NA} does \textbf{not} work. Always use
  \texttt{is.na()}
\end{itemize}
\end{frame}

\hypertarget{sec-factors}{%
\section{Factors}\label{sec-factors}}

\begin{frame}[fragile]{What Is a Factor?}
\protect\hypertarget{what-is-a-factor}{}
A \textbf{factor} stores a categorical variable with a fixed set of
\textbf{levels}.

\begin{Shaded}
\begin{Highlighting}[]
\NormalTok{region }\OtherTok{\textless{}{-}} \FunctionTok{c}\NormalTok{(}\StringTok{"South"}\NormalTok{, }\StringTok{"North"}\NormalTok{, }\StringTok{"West"}\NormalTok{, }\StringTok{"South"}\NormalTok{, }\StringTok{"North"}\NormalTok{, }\StringTok{"South"}\NormalTok{)}
\NormalTok{region\_f }\OtherTok{\textless{}{-}} \FunctionTok{factor}\NormalTok{(region)}

\NormalTok{region\_f}
\end{Highlighting}
\end{Shaded}

\begingroup\tiny\color{darkgray}

\begin{verbatim}
[1] South North West  South North South
Levels: North South West
\end{verbatim}

\endgroup

\begin{Shaded}
\begin{Highlighting}[]
\FunctionTok{levels}\NormalTok{(region\_f)        }\CommentTok{\# sorted alphabetically by default}
\end{Highlighting}
\end{Shaded}

\begingroup\tiny\color{darkgray}

\begin{verbatim}
[1] "North" "South" "West" 
\end{verbatim}

\endgroup

\begin{Shaded}
\begin{Highlighting}[]
\FunctionTok{table}\NormalTok{(region\_f)}
\end{Highlighting}
\end{Shaded}

\begingroup\tiny\color{darkgray}

\begin{verbatim}
region_f
North South  West 
    2     3     1 
\end{verbatim}

\endgroup

\begin{itemize}
\tightlist
\item
  For numeric codes, apply the data dictionary with \texttt{levels} and
  \texttt{labels} (Section~\ref{sec-coded})
\end{itemize}
\end{frame}

\begin{frame}[fragile]{Changing Levels with \texttt{forcats}}
\protect\hypertarget{changing-levels-with-forcats}{}
\texttt{forcats} (part of the tidyverse): every function starts with
\texttt{fct\_}

\begin{Shaded}
\begin{Highlighting}[]
\FunctionTok{fct\_recode}\NormalTok{(region\_f, }\StringTok{"Southeast"} \OtherTok{=} \StringTok{"South"}\NormalTok{)     }\CommentTok{\# rename: new = old}
\end{Highlighting}
\end{Shaded}

\begingroup\tiny\color{darkgray}

\begin{verbatim}
[1] Southeast North     West      Southeast North     Southeast
Levels: North Southeast West
\end{verbatim}

\endgroup

\begin{Shaded}
\begin{Highlighting}[]
\FunctionTok{fct\_collapse}\NormalTok{(satisfaction\_f,                    }\CommentTok{\# merge levels}
             \AttributeTok{Disagree =} \FunctionTok{c}\NormalTok{(}\StringTok{"Strongly disagree"}\NormalTok{, }\StringTok{"Disagree"}\NormalTok{),}
             \AttributeTok{Neutral  =} \StringTok{"Neutral"}\NormalTok{,}
             \AttributeTok{Agree    =} \FunctionTok{c}\NormalTok{(}\StringTok{"Agree"}\NormalTok{, }\StringTok{"Strongly agree"}\NormalTok{))}
\end{Highlighting}
\end{Shaded}

\begingroup\tiny\color{darkgray}

\begin{verbatim}
[1] Agree    Agree    Disagree Agree    Neutral  Disagree Agree   
[8] Agree   
Levels: Disagree < Neutral < Agree
\end{verbatim}

\endgroup

\begin{Shaded}
\begin{Highlighting}[]
\FunctionTok{fct\_infreq}\NormalTok{(region\_f)                            }\CommentTok{\# order by frequency}
\end{Highlighting}
\end{Shaded}

\begingroup\tiny\color{darkgray}

\begin{verbatim}
[1] South North West  South North South
Levels: South North West
\end{verbatim}

\endgroup
\end{frame}

\begin{frame}{The Reference Level: Why It Matters}
\protect\hypertarget{sec-reference}{}
\begin{itemize}
\tightlist
\item
  The \textbf{first level} of a factor is its \textbf{reference level}
  (baseline)
\item
  In a regression, \textbf{every other level is compared to the
  reference}
\item
  So the reference level decides \textbf{what your coefficients mean}
  (see Section~\ref{sec-lm})
\item
  Alphabetical order is rarely the comparison you want
\end{itemize}

\begin{exampleblock}{Choose the reference on purpose}
\protect\hypertarget{choose-the-reference-on-purpose}{}
\begin{itemize}
\tightlist
\item
  The \textbf{control group} in an experiment
\item
  The \textbf{status quo}, e.g.~``no training program''
\item
  The \textbf{largest} or most typical group
\item
  Whatever group your \textbf{hypothesis} compares against
\end{itemize}
\end{exampleblock}
\end{frame}

\begin{frame}[fragile]{The Reference Level: How to Set It}
\protect\hypertarget{the-reference-level-how-to-set-it}{}
\begin{Shaded}
\begin{Highlighting}[]
\CommentTok{\# 1. Base R: move one level to the front}
\FunctionTok{relevel}\NormalTok{(region\_f, }\AttributeTok{ref =} \StringTok{"West"}\NormalTok{) }\SpecialCharTok{|\textgreater{}} \FunctionTok{levels}\NormalTok{()}
\end{Highlighting}
\end{Shaded}

\begingroup\tiny\color{darkgray}

\begin{verbatim}
[1] "West"  "North" "South"
\end{verbatim}

\endgroup

\begin{Shaded}
\begin{Highlighting}[]
\CommentTok{\# 2. Base R: set the full order yourself}
\FunctionTok{factor}\NormalTok{(region, }\AttributeTok{levels =} \FunctionTok{c}\NormalTok{(}\StringTok{"West"}\NormalTok{, }\StringTok{"North"}\NormalTok{, }\StringTok{"South"}\NormalTok{)) }\SpecialCharTok{|\textgreater{}} \FunctionTok{levels}\NormalTok{()}
\end{Highlighting}
\end{Shaded}

\begingroup\tiny\color{darkgray}

\begin{verbatim}
[1] "West"  "North" "South"
\end{verbatim}

\endgroup

\begin{Shaded}
\begin{Highlighting}[]
\CommentTok{\# 3. forcats: move one or more levels to the front}
\FunctionTok{fct\_relevel}\NormalTok{(region\_f, }\StringTok{"West"}\NormalTok{) }\SpecialCharTok{|\textgreater{}} \FunctionTok{levels}\NormalTok{()}
\end{Highlighting}
\end{Shaded}

\begingroup\tiny\color{darkgray}

\begin{verbatim}
[1] "West"  "North" "South"
\end{verbatim}

\endgroup
\end{frame}

\begin{frame}[fragile]{A Common Trap}
\protect\hypertarget{a-common-trap}{}
\texttt{as.numeric()} on a factor returns the \textbf{internal codes},
not the original values:

\begin{Shaded}
\begin{Highlighting}[]
\NormalTok{year\_f }\OtherTok{\textless{}{-}} \FunctionTok{factor}\NormalTok{(}\FunctionTok{c}\NormalTok{(}\DecValTok{2022}\NormalTok{, }\DecValTok{2024}\NormalTok{, }\DecValTok{2023}\NormalTok{, }\DecValTok{2024}\NormalTok{))}

\FunctionTok{as.numeric}\NormalTok{(year\_f)                 }\CommentTok{\# 1 3 2 3 {-}{-} wrong!}
\end{Highlighting}
\end{Shaded}

\begingroup\tiny\color{darkgray}

\begin{verbatim}
[1] 1 3 2 3
\end{verbatim}

\endgroup

\begin{Shaded}
\begin{Highlighting}[]
\FunctionTok{as.numeric}\NormalTok{(}\FunctionTok{as.character}\NormalTok{(year\_f))   }\CommentTok{\# 2022 2024 2023 2024 {-}{-} right}
\end{Highlighting}
\end{Shaded}

\begingroup\tiny\color{darkgray}

\begin{verbatim}
[1] 2022 2024 2023 2024
\end{verbatim}

\endgroup
\end{frame}

\begin{frame}[fragile]{Task 1: Apply a Data Dictionary}
\protect\hypertarget{task-1-apply-a-data-dictionary}{}
\begin{alertblock}{Task}
\protect\hypertarget{task}{}
A customer file stores loyalty tier as codes: 1 = Bronze, 2 = Silver, 3
= Gold.

\begin{enumerate}
\tightlist
\item
  Convert \texttt{tier} to a factor with the correct labels
\item
  Count how many customers are in each tier using \texttt{table()}
\item
  Set \texttt{"Silver"} as the reference level and check with
  \texttt{levels()}
\item
  Use \texttt{fct\_collapse()} to combine Silver and Gold into
  \texttt{"Premium"}
\end{enumerate}

\begin{Shaded}
\begin{Highlighting}[]
\NormalTok{tier }\OtherTok{\textless{}{-}} \FunctionTok{c}\NormalTok{(}\DecValTok{3}\NormalTok{, }\DecValTok{1}\NormalTok{, }\DecValTok{2}\NormalTok{, }\DecValTok{3}\NormalTok{, }\DecValTok{1}\NormalTok{, }\DecValTok{3}\NormalTok{, }\DecValTok{2}\NormalTok{, }\DecValTok{1}\NormalTok{, }\DecValTok{1}\NormalTok{)}
\CommentTok{\# your code here}
\end{Highlighting}
\end{Shaded}
\end{alertblock}
\end{frame}

\begin{frame}[fragile]{Task 1: Solution}
\protect\hypertarget{task-1-solution}{}
\begin{Shaded}
\begin{Highlighting}[]
\NormalTok{tier }\OtherTok{\textless{}{-}} \FunctionTok{c}\NormalTok{(}\DecValTok{3}\NormalTok{, }\DecValTok{1}\NormalTok{, }\DecValTok{2}\NormalTok{, }\DecValTok{3}\NormalTok{, }\DecValTok{1}\NormalTok{, }\DecValTok{3}\NormalTok{, }\DecValTok{2}\NormalTok{, }\DecValTok{1}\NormalTok{, }\DecValTok{1}\NormalTok{)}

\NormalTok{tier\_f }\OtherTok{\textless{}{-}} \FunctionTok{factor}\NormalTok{(tier, }\AttributeTok{levels =} \DecValTok{1}\SpecialCharTok{:}\DecValTok{3}\NormalTok{, }\AttributeTok{labels =} \FunctionTok{c}\NormalTok{(}\StringTok{"Bronze"}\NormalTok{, }\StringTok{"Silver"}\NormalTok{, }\StringTok{"Gold"}\NormalTok{))}
\FunctionTok{table}\NormalTok{(tier\_f)}
\end{Highlighting}
\end{Shaded}

\begingroup\tiny\color{darkgray}

\begin{verbatim}
tier_f
Bronze Silver   Gold 
     4      2      3 
\end{verbatim}

\endgroup

\begin{Shaded}
\begin{Highlighting}[]
\NormalTok{tier\_f }\OtherTok{\textless{}{-}} \FunctionTok{fct\_relevel}\NormalTok{(tier\_f, }\StringTok{"Silver"}\NormalTok{)}
\FunctionTok{levels}\NormalTok{(tier\_f)}
\end{Highlighting}
\end{Shaded}

\begingroup\tiny\color{darkgray}

\begin{verbatim}
[1] "Silver" "Bronze" "Gold"  
\end{verbatim}

\endgroup

\begin{Shaded}
\begin{Highlighting}[]
\FunctionTok{fct\_collapse}\NormalTok{(tier\_f, }\AttributeTok{Premium =} \FunctionTok{c}\NormalTok{(}\StringTok{"Silver"}\NormalTok{, }\StringTok{"Gold"}\NormalTok{))}
\end{Highlighting}
\end{Shaded}

\begingroup\tiny\color{darkgray}

\begin{verbatim}
[1] Premium Bronze  Premium Premium Bronze  Premium Premium Bronze 
[9] Bronze 
Levels: Premium Bronze
\end{verbatim}

\endgroup
\end{frame}

\hypertarget{sec-structures}{%
\section{Data Structures}\label{sec-structures}}

\begin{frame}{R's Data Structures}
\protect\hypertarget{rs-data-structures}{}
\includegraphics[width=0.85\textwidth,height=\textheight]{images/data_structures.png}
\end{frame}

\begin{frame}{Comparing Data Structures}
\protect\hypertarget{comparing-data-structures}{}
\begin{longtable}[]{@{}llll@{}}
\toprule\noalign{}
Structure & Dimensions & Same type required? & Typical use \\
\midrule\noalign{}
\endhead
Vector & 1 & Yes & One variable (one column) \\
Matrix & 2 & Yes & Numeric computation \\
Array & 3+ & Yes & Stacks of matrices \\
Data frame & 2 & Within each column & Your dataset \\
List & 1 & No: anything & Model output, collections \\
\bottomrule\noalign{}
\end{longtable}
\end{frame}

\begin{frame}[fragile]{Vectors}
\protect\hypertarget{vectors}{}
A sequence of values of the \textbf{same type}. Every column in a
dataset is a vector.

\begin{Shaded}
\begin{Highlighting}[]
\NormalTok{sales }\OtherTok{\textless{}{-}} \FunctionTok{c}\NormalTok{(}\DecValTok{120}\NormalTok{, }\DecValTok{95}\NormalTok{, }\DecValTok{143}\NormalTok{, }\DecValTok{88}\NormalTok{, }\DecValTok{110}\NormalTok{)}
\FunctionTok{names}\NormalTok{(sales) }\OtherTok{\textless{}{-}} \FunctionTok{c}\NormalTok{(}\StringTok{"Mon"}\NormalTok{, }\StringTok{"Tue"}\NormalTok{, }\StringTok{"Wed"}\NormalTok{, }\StringTok{"Thu"}\NormalTok{, }\StringTok{"Fri"}\NormalTok{)}
\NormalTok{sales}
\end{Highlighting}
\end{Shaded}

\begingroup\tiny\color{darkgray}

\begin{verbatim}
Mon Tue Wed Thu Fri 
120  95 143  88 110 
\end{verbatim}

\endgroup

\begin{Shaded}
\begin{Highlighting}[]
\NormalTok{sales }\SpecialCharTok{*} \FloatTok{1.07}             \CommentTok{\# works on every element: no "drag down" as in Excel}
\end{Highlighting}
\end{Shaded}

\begingroup\tiny\color{darkgray}

\begin{verbatim}
   Mon    Tue    Wed    Thu    Fri 
128.40 101.65 153.01  94.16 117.70 
\end{verbatim}

\endgroup

\begin{Shaded}
\begin{Highlighting}[]
\FunctionTok{c}\NormalTok{(}\DecValTok{1}\NormalTok{, }\ConstantTok{TRUE}\NormalTok{, }\StringTok{"Sales"}\NormalTok{)      }\CommentTok{\# mixed types: everything becomes character}
\end{Highlighting}
\end{Shaded}

\begingroup\tiny\color{darkgray}

\begin{verbatim}
[1] "1"     "TRUE"  "Sales"
\end{verbatim}

\endgroup
\end{frame}

\begin{frame}{Subsetting a Vector}
\protect\hypertarget{sec-subset}{}
\includegraphics[width=0.85\textwidth,height=\textheight]{images/vector_subsetting.png}
\end{frame}

\begin{frame}[fragile]{Keep or Drop One Element}
\protect\hypertarget{keep-or-drop-one-element}{}
\begin{Shaded}
\begin{Highlighting}[]
\NormalTok{X }\OtherTok{\textless{}{-}} \FunctionTok{c}\NormalTok{(}\DecValTok{10}\NormalTok{, }\DecValTok{20}\NormalTok{, }\DecValTok{30}\NormalTok{, }\ConstantTok{NA}\NormalTok{, }\DecValTok{50}\NormalTok{)}
\FunctionTok{names}\NormalTok{(X) }\OtherTok{\textless{}{-}} \FunctionTok{c}\NormalTok{(}\StringTok{"a"}\NormalTok{, }\StringTok{"b"}\NormalTok{, }\StringTok{"c"}\NormalTok{, }\StringTok{"d"}\NormalTok{, }\StringTok{"e"}\NormalTok{)}
\end{Highlighting}
\end{Shaded}

\begin{columns}[T]
\begin{column}{0.5\textwidth}
\textbf{Keep} the third element

\begin{Shaded}
\begin{Highlighting}[]
\NormalTok{X[}\DecValTok{3}\NormalTok{]}
\end{Highlighting}
\end{Shaded}

\begingroup\tiny\color{darkgray}

\begin{verbatim}
 c 
30 
\end{verbatim}

\endgroup

\begin{Shaded}
\begin{Highlighting}[]
\NormalTok{X[}\StringTok{"c"}\NormalTok{]}
\end{Highlighting}
\end{Shaded}

\begingroup\tiny\color{darkgray}

\begin{verbatim}
 c 
30 
\end{verbatim}

\endgroup

\begin{Shaded}
\begin{Highlighting}[]
\NormalTok{X[}\FunctionTok{c}\NormalTok{(}\ConstantTok{FALSE}\NormalTok{, }\ConstantTok{FALSE}\NormalTok{, }\ConstantTok{TRUE}\NormalTok{, }\ConstantTok{FALSE}\NormalTok{, }\ConstantTok{FALSE}\NormalTok{)]}
\end{Highlighting}
\end{Shaded}

\begingroup\tiny\color{darkgray}

\begin{verbatim}
 c 
30 
\end{verbatim}

\endgroup
\end{column}

\begin{column}{0.5\textwidth}
\textbf{Drop} the third element

\begin{Shaded}
\begin{Highlighting}[]
\NormalTok{X[}\SpecialCharTok{{-}}\DecValTok{3}\NormalTok{]}
\end{Highlighting}
\end{Shaded}

\begingroup\tiny\color{darkgray}

\begin{verbatim}
 a  b  d  e 
10 20 NA 50 
\end{verbatim}

\endgroup

\begin{Shaded}
\begin{Highlighting}[]
\NormalTok{X[}\FunctionTok{names}\NormalTok{(X) }\SpecialCharTok{!=} \StringTok{"c"}\NormalTok{]}
\end{Highlighting}
\end{Shaded}

\begingroup\tiny\color{darkgray}

\begin{verbatim}
 a  b  d  e 
10 20 NA 50 
\end{verbatim}

\endgroup

\begin{Shaded}
\begin{Highlighting}[]
\NormalTok{X[}\SpecialCharTok{!}\FunctionTok{c}\NormalTok{(}\ConstantTok{FALSE}\NormalTok{, }\ConstantTok{FALSE}\NormalTok{, }\ConstantTok{TRUE}\NormalTok{, }\ConstantTok{FALSE}\NormalTok{, }\ConstantTok{FALSE}\NormalTok{)]}
\end{Highlighting}
\end{Shaded}

\begingroup\tiny\color{darkgray}

\begin{verbatim}
 a  b  d  e 
10 20 NA 50 
\end{verbatim}

\endgroup
\end{column}
\end{columns}
\end{frame}

\begin{frame}[fragile]{Drop Missing Values}
\protect\hypertarget{drop-missing-values}{}
All five return the elements that are not \texttt{NA}:

\begin{Shaded}
\begin{Highlighting}[]
\NormalTok{X[}\FunctionTok{which}\NormalTok{(}\SpecialCharTok{!}\FunctionTok{is.na}\NormalTok{(X))]          }\CommentTok{\# indexes of non{-}missing values}
\end{Highlighting}
\end{Shaded}

\begingroup\tiny\color{darkgray}

\begin{verbatim}
 a  b  c  e 
10 20 30 50 
\end{verbatim}

\endgroup

\begin{Shaded}
\begin{Highlighting}[]
\NormalTok{X[}\FunctionTok{names}\NormalTok{(X)[}\SpecialCharTok{!}\FunctionTok{is.na}\NormalTok{(X)]]       }\CommentTok{\# names of non{-}missing values}
\end{Highlighting}
\end{Shaded}

\begingroup\tiny\color{darkgray}

\begin{verbatim}
 a  b  c  e 
10 20 30 50 
\end{verbatim}

\endgroup

\begin{Shaded}
\begin{Highlighting}[]
\NormalTok{X[}\FunctionTok{is.na}\NormalTok{(X) }\SpecialCharTok{==} \ConstantTok{FALSE}\NormalTok{]         }\CommentTok{\# logical vector}
\end{Highlighting}
\end{Shaded}

\begingroup\tiny\color{darkgray}

\begin{verbatim}
 a  b  c  e 
10 20 30 50 
\end{verbatim}

\endgroup

\begin{Shaded}
\begin{Highlighting}[]
\NormalTok{X[}\FunctionTok{complete.cases}\NormalTok{(X)]         }\CommentTok{\# logical vector from complete.cases()}
\end{Highlighting}
\end{Shaded}

\begingroup\tiny\color{darkgray}

\begin{verbatim}
 a  b  c  e 
10 20 30 50 
\end{verbatim}

\endgroup

\begin{Shaded}
\begin{Highlighting}[]
\FunctionTok{na.omit}\NormalTok{(X)                   }\CommentTok{\# also records which were removed}
\end{Highlighting}
\end{Shaded}

\begingroup\tiny\color{darkgray}

\begin{verbatim}
 a  b  c  e 
10 20 30 50 
attr(,"na.action")
d 
4 
attr(,"class")
[1] "omit"
\end{verbatim}

\endgroup
\end{frame}

\begin{frame}[fragile]{Subset by Condition}
\protect\hypertarget{subset-by-condition}{}
The version you will use most: the logical vector is built from the data
itself.

\begin{Shaded}
\begin{Highlighting}[]
\NormalTok{sales[sales }\SpecialCharTok{\textgreater{}} \DecValTok{100}\NormalTok{]}
\end{Highlighting}
\end{Shaded}

\begingroup\tiny\color{darkgray}

\begin{verbatim}
Mon Wed Fri 
120 143 110 
\end{verbatim}

\endgroup

\begin{Shaded}
\begin{Highlighting}[]
\NormalTok{sales[sales }\SpecialCharTok{\textgreater{}} \DecValTok{100} \SpecialCharTok{\&}\NormalTok{ sales }\SpecialCharTok{\textless{}} \DecValTok{130}\NormalTok{]}
\end{Highlighting}
\end{Shaded}

\begingroup\tiny\color{darkgray}

\begin{verbatim}
Mon Fri 
120 110 
\end{verbatim}

\endgroup

\begin{exampleblock}{\texttt{TRUE}/\texttt{FALSE}
vs.~\texttt{T}/\texttt{F}}
\protect\hypertarget{truefalse-vs.-tf}{}
\texttt{T} and \texttt{F} work as shortcuts, but they are ordinary
variable names that can be overwritten
(e.g.~\texttt{T\ \textless{}-\ 0}). Spell them out in code you share.
\end{exampleblock}
\end{frame}

\begin{frame}[fragile]{Matrices and Arrays}
\protect\hypertarget{matrices-and-arrays}{}
\begin{columns}[T]
\begin{column}{0.5\textwidth}
\textbf{Matrix}: 2 dimensions, one type

\begin{Shaded}
\begin{Highlighting}[]
\NormalTok{m }\OtherTok{\textless{}{-}} \FunctionTok{matrix}\NormalTok{(}\DecValTok{1}\SpecialCharTok{:}\DecValTok{6}\NormalTok{, }\AttributeTok{nrow =} \DecValTok{2}\NormalTok{)}
\NormalTok{m}
\end{Highlighting}
\end{Shaded}

\begingroup\tiny\color{darkgray}

\begin{verbatim}
     [,1] [,2] [,3]
[1,]    1    3    5
[2,]    2    4    6
\end{verbatim}

\endgroup

\begin{Shaded}
\begin{Highlighting}[]
\NormalTok{m[}\DecValTok{2}\NormalTok{, }\DecValTok{3}\NormalTok{]     }\CommentTok{\# row 2, column 3}
\end{Highlighting}
\end{Shaded}

\begingroup\tiny\color{darkgray}

\begin{verbatim}
[1] 6
\end{verbatim}

\endgroup

\begin{Shaded}
\begin{Highlighting}[]
\NormalTok{m[}\DecValTok{1}\NormalTok{, ]      }\CommentTok{\# all of row 1}
\end{Highlighting}
\end{Shaded}

\begingroup\tiny\color{darkgray}

\begin{verbatim}
[1] 1 3 5
\end{verbatim}

\endgroup
\end{column}

\begin{column}{0.5\textwidth}
\textbf{Array}: 3+ dimensions, one type

\begin{Shaded}
\begin{Highlighting}[]
\NormalTok{a }\OtherTok{\textless{}{-}} \FunctionTok{array}\NormalTok{(}\DecValTok{1}\SpecialCharTok{:}\DecValTok{24}\NormalTok{, }\AttributeTok{dim =} \FunctionTok{c}\NormalTok{(}\DecValTok{2}\NormalTok{, }\DecValTok{3}\NormalTok{, }\DecValTok{4}\NormalTok{))}
\FunctionTok{dim}\NormalTok{(a)}
\end{Highlighting}
\end{Shaded}

\begingroup\tiny\color{darkgray}

\begin{verbatim}
[1] 2 3 4
\end{verbatim}

\endgroup

\begin{Shaded}
\begin{Highlighting}[]
\NormalTok{a[, , }\DecValTok{1}\NormalTok{]     }\CommentTok{\# first layer}
\end{Highlighting}
\end{Shaded}

\begingroup\tiny\color{darkgray}

\begin{verbatim}
     [,1] [,2] [,3]
[1,]    1    3    5
[2,]    2    4    6
\end{verbatim}

\endgroup

\begin{Shaded}
\begin{Highlighting}[]
\NormalTok{a[}\DecValTok{1}\NormalTok{, }\DecValTok{2}\NormalTok{, }\DecValTok{3}\NormalTok{]   }\CommentTok{\# row 1, col 2, layer 3}
\end{Highlighting}
\end{Shaded}

\begingroup\tiny\color{darkgray}

\begin{verbatim}
[1] 15
\end{verbatim}

\endgroup
\end{column}
\end{columns}
\end{frame}

\begin{frame}[fragile]{Data Frames}
\protect\hypertarget{sec-dataframes}{}
Looks like a matrix, but is a set of \textbf{equal-length vectors}: one
type per column.

\begin{Shaded}
\begin{Highlighting}[]
\NormalTok{employees }\OtherTok{\textless{}{-}} \FunctionTok{data.frame}\NormalTok{(}
  \AttributeTok{id           =} \DecValTok{1}\SpecialCharTok{:}\DecValTok{6}\NormalTok{,}
  \AttributeTok{department   =} \FunctionTok{c}\NormalTok{(}\StringTok{"Sales"}\NormalTok{, }\StringTok{"HR"}\NormalTok{, }\StringTok{"IT"}\NormalTok{, }\StringTok{"Sales"}\NormalTok{, }\StringTok{"IT"}\NormalTok{, }\StringTok{"HR"}\NormalTok{),}
  \AttributeTok{tenure\_years =} \FunctionTok{c}\NormalTok{(}\FloatTok{3.5}\NormalTok{, }\DecValTok{10}\NormalTok{, }\FloatTok{1.2}\NormalTok{, }\DecValTok{7}\NormalTok{, }\ConstantTok{NA}\NormalTok{, }\FloatTok{4.8}\NormalTok{),}
  \AttributeTok{remote       =} \FunctionTok{c}\NormalTok{(}\ConstantTok{TRUE}\NormalTok{, }\ConstantTok{FALSE}\NormalTok{, }\ConstantTok{TRUE}\NormalTok{, }\ConstantTok{TRUE}\NormalTok{, }\ConstantTok{FALSE}\NormalTok{, }\ConstantTok{FALSE}\NormalTok{),}
  \AttributeTok{hire\_date    =} \FunctionTok{ymd}\NormalTok{(}\FunctionTok{c}\NormalTok{(}\StringTok{"2023{-}05{-}01"}\NormalTok{, }\StringTok{"2016{-}11{-}14"}\NormalTok{, }\StringTok{"2025{-}08{-}20"}\NormalTok{,}
                       \StringTok{"2019{-}11{-}04"}\NormalTok{, }\StringTok{"2026{-}01{-}12"}\NormalTok{, }\StringTok{"2022{-}01{-}03"}\NormalTok{))}
\NormalTok{)}
\FunctionTok{str}\NormalTok{(employees)}
\end{Highlighting}
\end{Shaded}

\begingroup\tiny\color{darkgray}

\begin{verbatim}
'data.frame':   6 obs. of  5 variables:
 $ id          : int  1 2 3 4 5 6
 $ department  : chr  "Sales" "HR" "IT" "Sales" ...
 $ tenure_years: num  3.5 10 1.2 7 NA 4.8
 $ remote      : logi  TRUE FALSE TRUE TRUE FALSE FALSE
 $ hire_date   : Date, format: "2023-05-01" ...
\end{verbatim}

\endgroup
\end{frame}

\begin{frame}[fragile]{Subsetting a Data Frame}
\protect\hypertarget{subsetting-a-data-frame}{}
\texttt{data{[}rows,\ columns{]}}: a blank side means ``all''

\begin{Shaded}
\begin{Highlighting}[]
\NormalTok{employees}\SpecialCharTok{$}\NormalTok{department                    }\CommentTok{\# one column as a vector}
\end{Highlighting}
\end{Shaded}

\begingroup\tiny\color{darkgray}

\begin{verbatim}
[1] "Sales" "HR"    "IT"    "Sales" "IT"    "HR"   
\end{verbatim}

\endgroup

\begin{Shaded}
\begin{Highlighting}[]
\NormalTok{employees[employees}\SpecialCharTok{$}\NormalTok{remote, ]           }\CommentTok{\# rows where remote is TRUE}
\end{Highlighting}
\end{Shaded}

\begingroup\tiny\color{darkgray}

\begin{verbatim}
  id department tenure_years remote  hire_date
1  1      Sales          3.5   TRUE 2023-05-01
3  3         IT          1.2   TRUE 2025-08-20
4  4      Sales          7.0   TRUE 2019-11-04
\end{verbatim}

\endgroup

\begin{Shaded}
\begin{Highlighting}[]
\NormalTok{employees[, }\FunctionTok{c}\NormalTok{(}\StringTok{"id"}\NormalTok{, }\StringTok{"tenure\_years"}\NormalTok{)]    }\CommentTok{\# selected columns}
\end{Highlighting}
\end{Shaded}

\begingroup\tiny\color{darkgray}

\begin{verbatim}
  id tenure_years
1  1          3.5
2  2         10.0
3  3          1.2
4  4          7.0
5  5           NA
6  6          4.8
\end{verbatim}

\endgroup
\end{frame}

\begin{frame}[fragile]{Modifying a Data Frame}
\protect\hypertarget{modifying-a-data-frame}{}
\begin{Shaded}
\begin{Highlighting}[]
\NormalTok{employees}\SpecialCharTok{$}\NormalTok{department }\OtherTok{\textless{}{-}} \FunctionTok{factor}\NormalTok{(employees}\SpecialCharTok{$}\NormalTok{department)  }\CommentTok{\# convert to factor}
\NormalTok{employees}\SpecialCharTok{$}\NormalTok{senior }\OtherTok{\textless{}{-}}\NormalTok{ employees}\SpecialCharTok{$}\NormalTok{tenure\_years }\SpecialCharTok{\textgreater{}=} \DecValTok{5}       \CommentTok{\# add a new column}
\end{Highlighting}
\end{Shaded}

Preview of \texttt{dplyr} (Sessions 5 and 6): the same ideas, read like
sentences

\begin{Shaded}
\begin{Highlighting}[]
\NormalTok{employees }\SpecialCharTok{|\textgreater{}}
  \FunctionTok{filter}\NormalTok{(remote) }\SpecialCharTok{|\textgreater{}}
  \FunctionTok{select}\NormalTok{(id, department, senior)}
\end{Highlighting}
\end{Shaded}

\begingroup\tiny\color{darkgray}

\begin{verbatim}
  id department senior
1  1      Sales  FALSE
2  3         IT  FALSE
3  4      Sales   TRUE
\end{verbatim}

\endgroup

\begin{itemize}
\tightlist
\item
  Tidyverse functions often return a \textbf{tibble}: a modern data
  frame that prints more neatly
\end{itemize}
\end{frame}

\begin{frame}[fragile]{Lists}
\protect\hypertarget{lists}{}
One-dimensional like a vector, but each element can be
\textbf{anything}.

\begin{Shaded}
\begin{Highlighting}[]
\NormalTok{firm }\OtherTok{\textless{}{-}} \FunctionTok{list}\NormalTok{(}
  \AttributeTok{name      =} \StringTok{"Acme Corp"}\NormalTok{,}
  \AttributeTok{founded   =}\NormalTok{ 1998L,}
  \AttributeTok{divisions =} \FunctionTok{c}\NormalTok{(}\StringTok{"Retail"}\NormalTok{, }\StringTok{"Wholesale"}\NormalTok{, }\StringTok{"Online"}\NormalTok{),}
  \AttributeTok{staff     =}\NormalTok{ employees}
\NormalTok{)}

\FunctionTok{str}\NormalTok{(firm, }\AttributeTok{max.level =} \DecValTok{1}\NormalTok{)}
\end{Highlighting}
\end{Shaded}

\begingroup\tiny\color{darkgray}

\begin{verbatim}
List of 4
 $ name     : chr "Acme Corp"
 $ founded  : int 1998
 $ divisions: chr [1:3] "Retail" "Wholesale" "Online"
 $ staff    :'data.frame':  6 obs. of  6 variables:
\end{verbatim}

\endgroup

\begin{Shaded}
\begin{Highlighting}[]
\NormalTok{firm}\SpecialCharTok{$}\NormalTok{divisions; firm[[}\StringTok{"founded"}\NormalTok{]]}
\end{Highlighting}
\end{Shaded}

\begingroup\tiny\color{darkgray}

\begin{verbatim}
[1] "Retail"    "Wholesale" "Online"   
\end{verbatim}

\endgroup

\begingroup\tiny\color{darkgray}

\begin{verbatim}
[1] 1998
\end{verbatim}

\endgroup
\end{frame}

\hypertarget{sec-lm}{%
\section{Factors in a Regression Model}\label{sec-lm}}

\begin{frame}[fragile]{The \texttt{Salaries} Data}
\protect\hypertarget{the-salaries-data}{}
\begin{itemize}
\tightlist
\item
  2008--09 nine-month salaries of 397 professors at a US university
  (package \texttt{carData})
\item
  You know regression already. Here we only \textbf{read the
  coefficients}
\end{itemize}

\begin{Shaded}
\begin{Highlighting}[]
\FunctionTok{str}\NormalTok{(Salaries)}
\end{Highlighting}
\end{Shaded}

\begingroup\tiny\color{darkgray}

\begin{verbatim}
'data.frame':   397 obs. of  6 variables:
 $ rank         : Factor w/ 3 levels "AsstProf","AssocProf",..: 3 3 1 3 3 2 3 3 3 3 ...
 $ discipline   : Factor w/ 2 levels "A","B": 2 2 2 2 2 2 2 2 2 2 ...
 $ yrs.since.phd: int  19 20 4 45 40 6 30 45 21 18 ...
 $ yrs.service  : int  18 16 3 39 41 6 23 45 20 18 ...
 $ sex          : Factor w/ 2 levels "Female","Male": 2 2 2 2 2 2 2 2 2 1 ...
 $ salary       : int  139750 173200 79750 115000 141500 97000 175000 147765 119250 129000 ...
\end{verbatim}

\endgroup
\end{frame}

\begin{frame}[fragile]{Average Salary by Rank}
\protect\hypertarget{average-salary-by-rank}{}
\begin{Shaded}
\begin{Highlighting}[]
\NormalTok{Salaries }\SpecialCharTok{|\textgreater{}}
  \FunctionTok{group\_by}\NormalTok{(rank) }\SpecialCharTok{|\textgreater{}}
  \FunctionTok{summarise}\NormalTok{(}\AttributeTok{n =} \FunctionTok{n}\NormalTok{(), }\AttributeTok{mean\_salary =} \FunctionTok{mean}\NormalTok{(salary))}
\end{Highlighting}
\end{Shaded}

\begingroup\tiny\color{darkgray}

\begin{verbatim}
# A tibble: 3 x 3
  rank          n mean_salary
  <fct>     <int>       <dbl>
1 AsstProf     67      80776.
2 AssocProf    64      93876.
3 Prof        266     126772.
\end{verbatim}

\endgroup

\begin{Shaded}
\begin{Highlighting}[]
\NormalTok{group\_means }\OtherTok{\textless{}{-}} \FunctionTok{tapply}\NormalTok{(Salaries}\SpecialCharTok{$}\NormalTok{salary, Salaries}\SpecialCharTok{$}\NormalTok{rank, mean)   }\CommentTok{\# base R version}
\end{Highlighting}
\end{Shaded}
\end{frame}

\begin{frame}[fragile]{Reference Coding (R's Default)}
\protect\hypertarget{sec-ref-coding}{}
\begin{Shaded}
\begin{Highlighting}[]
\NormalTok{m\_ref }\OtherTok{\textless{}{-}} \FunctionTok{lm}\NormalTok{(salary }\SpecialCharTok{\textasciitilde{}}\NormalTok{ rank, }\AttributeTok{data =}\NormalTok{ Salaries)}
\FunctionTok{summary}\NormalTok{(m\_ref)}
\end{Highlighting}
\end{Shaded}

\begingroup\tiny\color{darkgray}

\begin{verbatim}

Call:
lm(formula = salary ~ rank, data = Salaries)

Residuals:
   Min     1Q Median     3Q    Max 
-68972 -16376  -1580  11755 104773 

Coefficients:
              Estimate Std. Error t value Pr(>|t|)    
(Intercept)      80776       2887  27.976  < 2e-16 ***
rankAssocProf    13100       4131   3.171  0.00164 ** 
rankProf         45996       3230  14.238  < 2e-16 ***
---
Signif. codes:  0 '***' 0.001 '**' 0.01 '*' 0.05 '.' 0.1 ' ' 1

Residual standard error: 23630 on 394 degrees of freedom
Multiple R-squared:  0.3943,    Adjusted R-squared:  0.3912 
F-statistic: 128.2 on 2 and 394 DF,  p-value: < 2.2e-16
\end{verbatim}

\endgroup
\end{frame}

\begin{frame}[fragile]{What R Built Behind the Scenes}
\protect\hypertarget{what-r-built-behind-the-scenes}{}
R turned \texttt{rank} into \textbf{two} dummy (0/1) columns: one per
level except the reference.

\begin{Shaded}
\begin{Highlighting}[]
\FunctionTok{model.matrix}\NormalTok{(m\_ref) }\SpecialCharTok{|\textgreater{}} \FunctionTok{head}\NormalTok{()}
\end{Highlighting}
\end{Shaded}

\begingroup\tiny\color{darkgray}

\begin{verbatim}
  (Intercept) rankAssocProf rankProf
1           1             0        1
2           1             0        1
3           1             0        0
4           1             0        1
5           1             0        1
6           1             1        0
\end{verbatim}

\endgroup
\end{frame}

\begin{frame}[fragile]{Reading the Coefficients}
\protect\hypertarget{reading-the-coefficients}{}
\begin{itemize}
\tightlist
\item
  \textbf{\texttt{(Intercept)}} = \$80,776: the average salary of the
  reference level, \textbf{AsstProf}
\item
  \textbf{rankAssocProf} = \$13,100: how much more AssocProf professors
  earn \textbf{than AsstProf}
\item
  \textbf{rankProf} = \$45,996: how much more Prof professors earn
  \textbf{than AsstProf}
\end{itemize}

\begin{exampleblock}{Key idea}
\protect\hypertarget{key-idea}{}
Each coefficient is a \textbf{difference from the reference level}. Its
p-value tests whether that difference is zero.
\end{exampleblock}
\end{frame}

\begin{frame}[fragile]{Changing the Reference Level}
\protect\hypertarget{changing-the-reference-level}{}
Research question: how much less do junior faculty earn than
\textbf{full professors}?

\begin{Shaded}
\begin{Highlighting}[]
\NormalTok{Salaries}\SpecialCharTok{$}\NormalTok{rank\_prof }\OtherTok{\textless{}{-}} \FunctionTok{fct\_relevel}\NormalTok{(Salaries}\SpecialCharTok{$}\NormalTok{rank, }\StringTok{"Prof"}\NormalTok{)}
\NormalTok{m\_ref2 }\OtherTok{\textless{}{-}} \FunctionTok{lm}\NormalTok{(salary }\SpecialCharTok{\textasciitilde{}}\NormalTok{ rank\_prof, }\AttributeTok{data =}\NormalTok{ Salaries)}
\FunctionTok{coef}\NormalTok{(m\_ref2)}
\end{Highlighting}
\end{Shaded}

\begingroup\tiny\color{darkgray}

\begin{verbatim}
       (Intercept)  rank_profAsstProf rank_profAssocProf 
         126772.11          -45996.12          -32895.67 
\end{verbatim}

\endgroup

\begin{Shaded}
\begin{Highlighting}[]
\FunctionTok{all.equal}\NormalTok{(}\FunctionTok{fitted}\NormalTok{(m\_ref), }\FunctionTok{fitted}\NormalTok{(m\_ref2))   }\CommentTok{\# identical predictions}
\end{Highlighting}
\end{Shaded}

\begingroup\tiny\color{darkgray}

\begin{verbatim}
[1] TRUE
\end{verbatim}

\endgroup

\begin{itemize}
\tightlist
\item
  Every coefficient changed, but the \textbf{model did not}
\item
  The reference level changes \emph{the question the coefficients
  answer}, not the fit
\end{itemize}
\end{frame}

\begin{frame}[fragile]{One-Hot Coding (No Intercept)}
\protect\hypertarget{sec-onehot}{}
One 0/1 column for \textbf{every} level. Remove the intercept with
\texttt{0\ +}:

\begin{Shaded}
\begin{Highlighting}[]
\NormalTok{m\_onehot }\OtherTok{\textless{}{-}} \FunctionTok{lm}\NormalTok{(salary }\SpecialCharTok{\textasciitilde{}} \DecValTok{0} \SpecialCharTok{+}\NormalTok{ rank, }\AttributeTok{data =}\NormalTok{ Salaries)}
\FunctionTok{coef}\NormalTok{(m\_onehot)}
\end{Highlighting}
\end{Shaded}

\begingroup\tiny\color{darkgray}

\begin{verbatim}
 rankAsstProf rankAssocProf      rankProf 
     80775.99      93876.44     126772.11 
\end{verbatim}

\endgroup

\begin{Shaded}
\begin{Highlighting}[]
\FunctionTok{model.matrix}\NormalTok{(m\_onehot) }\SpecialCharTok{|\textgreater{}} \FunctionTok{head}\NormalTok{(}\DecValTok{3}\NormalTok{)}
\end{Highlighting}
\end{Shaded}

\begingroup\tiny\color{darkgray}

\begin{verbatim}
  rankAsstProf rankAssocProf rankProf
1            0             0        1
2            0             0        1
3            1             0        0
\end{verbatim}

\endgroup

\begin{itemize}
\tightlist
\item
  No intercept; each coefficient is simply \textbf{the average salary of
  that rank}
\end{itemize}
\end{frame}

\begin{frame}[fragile]{Same Information, Different Numbers}
\protect\hypertarget{same-information-different-numbers}{}
\begin{Shaded}
\begin{Highlighting}[]
\FunctionTok{data.frame}\NormalTok{(}\AttributeTok{group\_mean       =}\NormalTok{ group\_means,}
           \AttributeTok{reference\_coding =} \FunctionTok{coef}\NormalTok{(m\_ref),}
           \AttributeTok{one\_hot          =} \FunctionTok{coef}\NormalTok{(m\_onehot))}
\end{Highlighting}
\end{Shaded}

\begingroup\tiny\color{darkgray}

\begin{verbatim}
          group_mean reference_coding   one_hot
AsstProf    80775.99         80775.99  80775.99
AssocProf   93876.44         13100.45  93876.44
Prof       126772.11         45996.12 126772.11
\end{verbatim}

\endgroup

\begin{Shaded}
\begin{Highlighting}[]
\FunctionTok{all.equal}\NormalTok{(}\FunctionTok{fitted}\NormalTok{(m\_ref), }\FunctionTok{fitted}\NormalTok{(m\_onehot))}
\end{Highlighting}
\end{Shaded}

\begingroup\tiny\color{darkgray}

\begin{verbatim}
[1] TRUE
\end{verbatim}

\endgroup

\begin{itemize}
\tightlist
\item
  Reference coding: AssocProf average = \$80,776 + \$13,100 = \$93,876
\end{itemize}
\end{frame}

\begin{frame}[fragile]{Reference Coding vs.~One-Hot Coding}
\protect\hypertarget{reference-coding-vs.-one-hot-coding}{}
\begin{longtable}[]{@{}
  >{\raggedright\arraybackslash}p{(\columnwidth - 4\tabcolsep) * \real{0.3333}}
  >{\raggedright\arraybackslash}p{(\columnwidth - 4\tabcolsep) * \real{0.3333}}
  >{\raggedright\arraybackslash}p{(\columnwidth - 4\tabcolsep) * \real{0.3333}}@{}}
\toprule\noalign{}
\begin{minipage}[b]{\linewidth}\raggedright
\end{minipage} & \begin{minipage}[b]{\linewidth}\raggedright
Reference coding
\end{minipage} & \begin{minipage}[b]{\linewidth}\raggedright
One-hot coding
\end{minipage} \\
\midrule\noalign{}
\endhead
Formula & \texttt{salary\ \textasciitilde{}\ rank} &
\texttt{salary\ \textasciitilde{}\ 0\ +\ rank} \\
Dummy columns & One per level, minus the reference & One per level \\
Intercept & Mean of the reference level & None \\
Coefficients & Differences from the reference & Group means \\
p-values test & Different from the reference? & Is the mean zero?
(rarely useful) \\
Fitted values & Same & Same \\
\bottomrule\noalign{}
\end{longtable}

\begin{alertblock}{Don't compare R² across these models}
\protect\hypertarget{dont-compare-ruxb2-across-these-models}{}
Without an intercept, \texttt{summary()} computes R² differently, so it
looks much higher even though the fit is identical.
\end{alertblock}
\end{frame}

\begin{frame}[fragile]{Ordered Factors in a Regression}
\protect\hypertarget{sec-ordered}{}
\begin{itemize}
\tightlist
\item
  \texttt{diamonds} (from \texttt{ggplot2}): prices and attributes of
  about 54,000 diamonds
\item
  Regressing price on product attributes is a \textbf{hedonic pricing
  model}, a standard method in marketing and economics
\item
  \texttt{cut} is an \textbf{ordered} factor:
\end{itemize}

\begin{Shaded}
\begin{Highlighting}[]
\FunctionTok{class}\NormalTok{(diamonds}\SpecialCharTok{$}\NormalTok{cut)}
\end{Highlighting}
\end{Shaded}

\begingroup\tiny\color{darkgray}

\begin{verbatim}
[1] "ordered" "factor" 
\end{verbatim}

\endgroup

\begin{Shaded}
\begin{Highlighting}[]
\FunctionTok{levels}\NormalTok{(diamonds}\SpecialCharTok{$}\NormalTok{cut)}
\end{Highlighting}
\end{Shaded}

\begingroup\tiny\color{darkgray}

\begin{verbatim}
[1] "Fair"      "Good"      "Very Good" "Premium"   "Ideal"    
\end{verbatim}

\endgroup
\end{frame}

\begin{frame}[fragile]{Polynomial Contrasts}
\protect\hypertarget{polynomial-contrasts}{}
\begin{Shaded}
\begin{Highlighting}[]
\NormalTok{m\_ordered }\OtherTok{\textless{}{-}} \FunctionTok{lm}\NormalTok{(price }\SpecialCharTok{\textasciitilde{}}\NormalTok{ cut, }\AttributeTok{data =}\NormalTok{ diamonds)}
\FunctionTok{coef}\NormalTok{(m\_ordered)}
\end{Highlighting}
\end{Shaded}

\begingroup\tiny\color{darkgray}

\begin{verbatim}
(Intercept)       cut.L       cut.Q       cut.C       cut^4 
  4062.2364   -362.7254   -225.5798   -699.4965   -280.3564 
\end{verbatim}

\endgroup

\begin{itemize}
\tightlist
\item
  \texttt{cut.L}, \texttt{cut.Q}, \texttt{cut.C}, \texttt{cut\^{}4}:
  \textbf{linear, quadratic, cubic, and fourth-power trends} across the
  ordered levels
\item
  A valid coding, but \textbf{not} ``difference from the reference''
\end{itemize}
\end{frame}

\begin{frame}[fragile]{Back to Reference Coding}
\protect\hypertarget{back-to-reference-coding}{}
\begin{Shaded}
\begin{Highlighting}[]
\NormalTok{diamonds}\SpecialCharTok{$}\NormalTok{cut\_unordered }\OtherTok{\textless{}{-}} \FunctionTok{factor}\NormalTok{(diamonds}\SpecialCharTok{$}\NormalTok{cut, }\AttributeTok{ordered =} \ConstantTok{FALSE}\NormalTok{)}
\NormalTok{m\_unordered }\OtherTok{\textless{}{-}} \FunctionTok{lm}\NormalTok{(price }\SpecialCharTok{\textasciitilde{}}\NormalTok{ cut\_unordered, }\AttributeTok{data =}\NormalTok{ diamonds)}
\FunctionTok{coef}\NormalTok{(m\_unordered)}
\end{Highlighting}
\end{Shaded}

\begingroup\tiny\color{darkgray}

\begin{verbatim}
           (Intercept)      cut_unorderedGood cut_unorderedVery Good 
             4358.7578              -429.8933              -376.9979 
  cut_unorderedPremium     cut_unorderedIdeal 
              225.4999              -901.2158 
\end{verbatim}

\endgroup

\begin{Shaded}
\begin{Highlighting}[]
\FunctionTok{all.equal}\NormalTok{(}\FunctionTok{fitted}\NormalTok{(m\_ordered), }\FunctionTok{fitted}\NormalTok{(m\_unordered))   }\CommentTok{\# same predictions again}
\end{Highlighting}
\end{Shaded}

\begingroup\tiny\color{darkgray}

\begin{verbatim}
[1] TRUE
\end{verbatim}

\endgroup

\begin{alertblock}{Ordered factors and your survey data}
\protect\hypertarget{ordered-factors-and-your-survey-data}{}
Likert items stored as ordered factors (Section~\ref{sec-likert}) get
polynomial contrasts. Unless you want trend tests, convert them to
unordered factors or keep them numeric before modeling.
\end{alertblock}
\end{frame}

\begin{frame}[fragile]{A Puzzle for Later}
\protect\hypertarget{a-puzzle-for-later}{}
\begin{itemize}
\tightlist
\item
  On average, \texttt{Ideal} diamonds are \emph{cheaper} than
  \texttt{Fair} ones. Better cut, lower price?
\item
  Diamond size (\texttt{carat}) is missing from the model, and ideal-cut
  diamonds tend to be smaller
\item
  This is \textbf{omitted variable bias}, which your statistics course
  will cover
\end{itemize}
\end{frame}

\begin{frame}[fragile]{The Model Is a List}
\protect\hypertarget{the-model-is-a-list}{}
A fitted model is one of our data structures: a \textbf{list}.

\begin{Shaded}
\begin{Highlighting}[]
\FunctionTok{is.list}\NormalTok{(m\_ref)}
\end{Highlighting}
\end{Shaded}

\begingroup\tiny\color{darkgray}

\begin{verbatim}
[1] TRUE
\end{verbatim}

\endgroup

\begin{Shaded}
\begin{Highlighting}[]
\FunctionTok{names}\NormalTok{(m\_ref)}
\end{Highlighting}
\end{Shaded}

\begingroup\tiny\color{darkgray}

\begin{verbatim}
 [1] "coefficients"  "residuals"     "effects"       "rank"         
 [5] "fitted.values" "assign"        "qr"            "df.residual"  
 [9] "contrasts"     "xlevels"       "call"          "terms"        
[13] "model"        
\end{verbatim}

\endgroup

\begin{Shaded}
\begin{Highlighting}[]
\NormalTok{m\_ref}\SpecialCharTok{$}\NormalTok{coefficients}
\end{Highlighting}
\end{Shaded}

\begingroup\tiny\color{darkgray}

\begin{verbatim}
  (Intercept) rankAssocProf      rankProf 
     80775.99      13100.45      45996.12 
\end{verbatim}

\endgroup

\begin{itemize}
\tightlist
\item
  This is why \texttt{\$} works on model objects
\end{itemize}
\end{frame}

\hypertarget{sec-your-turn}{%
\section{Your Turn}\label{sec-your-turn}}

\begin{frame}[fragile]{Exercise}
\protect\hypertarget{exercise}{}
\begin{exampleblock}{Data dictionary}
\protect\hypertarget{data-dictionary}{}
\texttt{region}: 1 = East, 2 = South, 3 = West\\
\texttt{satisfaction}: 1 = Very dissatisfied, 2 = Dissatisfied, 3 =
Neutral, 4 = Satisfied, 5 = Very satisfied

\begin{enumerate}
\tightlist
\item
  Check every column's type with \texttt{str()}. Which are stored in a
  way that does not match their meaning?
\item
  Convert \texttt{region} to a factor using the dictionary, with
  \texttt{"South"} as the reference level
\item
  Convert \texttt{annual\_spend} to numeric. What happens to
  \texttt{"unknown"}? Count the missing values and compute the average
  spend
\item
  Create \texttt{satisfaction\_f}, an ordered factor version of
  \texttt{satisfaction}. Keep the numeric column too
\item
  Convert \texttt{signup\_date} to a Date and compute tenure in days as
  of \texttt{2026-11-04}
\item
  Keep only customers who renewed; show just \texttt{customer} and
  \texttt{annual\_spend}
\end{enumerate}
\end{exampleblock}
\end{frame}

\begin{frame}[fragile]{Exercise Data}
\protect\hypertarget{exercise-data}{}
\begin{Shaded}
\begin{Highlighting}[]
\NormalTok{survey }\OtherTok{\textless{}{-}} \FunctionTok{data.frame}\NormalTok{(}
  \AttributeTok{customer     =} \FunctionTok{c}\NormalTok{(}\StringTok{"C01"}\NormalTok{, }\StringTok{"C02"}\NormalTok{, }\StringTok{"C03"}\NormalTok{, }\StringTok{"C04"}\NormalTok{, }\StringTok{"C05"}\NormalTok{, }\StringTok{"C06"}\NormalTok{, }\StringTok{"C07"}\NormalTok{, }\StringTok{"C08"}\NormalTok{),}
  \AttributeTok{region       =} \FunctionTok{c}\NormalTok{(}\DecValTok{3}\NormalTok{, }\DecValTok{1}\NormalTok{, }\DecValTok{2}\NormalTok{, }\DecValTok{1}\NormalTok{, }\DecValTok{3}\NormalTok{, }\DecValTok{2}\NormalTok{, }\DecValTok{2}\NormalTok{, }\DecValTok{1}\NormalTok{),}
  \AttributeTok{annual\_spend =} \FunctionTok{c}\NormalTok{(}\StringTok{"1200"}\NormalTok{, }\StringTok{"850"}\NormalTok{, }\StringTok{"unknown"}\NormalTok{, }\StringTok{"2300"}\NormalTok{, }\StringTok{"1750"}\NormalTok{, }\StringTok{"640"}\NormalTok{, }\StringTok{"990"}\NormalTok{, }\StringTok{"1480"}\NormalTok{),}
  \AttributeTok{satisfaction =} \FunctionTok{c}\NormalTok{(}\DecValTok{5}\NormalTok{, }\DecValTok{2}\NormalTok{, }\DecValTok{3}\NormalTok{, }\DecValTok{4}\NormalTok{, }\DecValTok{4}\NormalTok{, }\DecValTok{1}\NormalTok{, }\DecValTok{3}\NormalTok{, }\DecValTok{5}\NormalTok{),}
  \AttributeTok{signup\_date  =} \FunctionTok{c}\NormalTok{(}\StringTok{"03/15/2022"}\NormalTok{, }\StringTok{"11/02/2024"}\NormalTok{, }\StringTok{"07/30/2019"}\NormalTok{, }\StringTok{"01/09/2025"}\NormalTok{,}
                   \StringTok{"06/21/2021"}\NormalTok{, }\StringTok{"09/14/2023"}\NormalTok{, }\StringTok{"12/01/2020"}\NormalTok{, }\StringTok{"04/18/2026"}\NormalTok{),}
  \AttributeTok{renewed      =} \FunctionTok{c}\NormalTok{(}\ConstantTok{TRUE}\NormalTok{, }\ConstantTok{FALSE}\NormalTok{, }\ConstantTok{TRUE}\NormalTok{, }\ConstantTok{TRUE}\NormalTok{, }\ConstantTok{FALSE}\NormalTok{, }\ConstantTok{TRUE}\NormalTok{, }\ConstantTok{FALSE}\NormalTok{, }\ConstantTok{TRUE}\NormalTok{)}
\NormalTok{)}
\end{Highlighting}
\end{Shaded}
\end{frame}

\begin{frame}[fragile]{Exercise: Solution (1--3)}
\protect\hypertarget{exercise-solution-13}{}
\begin{Shaded}
\begin{Highlighting}[]
\CommentTok{\# 1. region and satisfaction are codes stored as numbers;}
\CommentTok{\#    annual\_spend and signup\_date are stored as character}

\CommentTok{\# 2.}
\NormalTok{survey}\SpecialCharTok{$}\NormalTok{region }\OtherTok{\textless{}{-}} \FunctionTok{factor}\NormalTok{(survey}\SpecialCharTok{$}\NormalTok{region, }\AttributeTok{levels =} \DecValTok{1}\SpecialCharTok{:}\DecValTok{3}\NormalTok{,}
                        \AttributeTok{labels =} \FunctionTok{c}\NormalTok{(}\StringTok{"East"}\NormalTok{, }\StringTok{"South"}\NormalTok{, }\StringTok{"West"}\NormalTok{))}
\NormalTok{survey}\SpecialCharTok{$}\NormalTok{region }\OtherTok{\textless{}{-}} \FunctionTok{fct\_relevel}\NormalTok{(survey}\SpecialCharTok{$}\NormalTok{region, }\StringTok{"South"}\NormalTok{)}
\FunctionTok{levels}\NormalTok{(survey}\SpecialCharTok{$}\NormalTok{region)}
\end{Highlighting}
\end{Shaded}

\begingroup\tiny\color{darkgray}

\begin{verbatim}
[1] "South" "East"  "West" 
\end{verbatim}

\endgroup

\begin{Shaded}
\begin{Highlighting}[]
\CommentTok{\# 3. "unknown" becomes NA (with a warning)}
\NormalTok{survey}\SpecialCharTok{$}\NormalTok{annual\_spend }\OtherTok{\textless{}{-}} \FunctionTok{as.numeric}\NormalTok{(survey}\SpecialCharTok{$}\NormalTok{annual\_spend)}
\FunctionTok{sum}\NormalTok{(}\FunctionTok{is.na}\NormalTok{(survey}\SpecialCharTok{$}\NormalTok{annual\_spend))}
\end{Highlighting}
\end{Shaded}

\begingroup\tiny\color{darkgray}

\begin{verbatim}
[1] 1
\end{verbatim}

\endgroup

\begin{Shaded}
\begin{Highlighting}[]
\FunctionTok{mean}\NormalTok{(survey}\SpecialCharTok{$}\NormalTok{annual\_spend, }\AttributeTok{na.rm =} \ConstantTok{TRUE}\NormalTok{)}
\end{Highlighting}
\end{Shaded}

\begingroup\tiny\color{darkgray}

\begin{verbatim}
[1] 1315.714
\end{verbatim}

\endgroup
\end{frame}

\begin{frame}[fragile]{Exercise: Solution (4--6)}
\protect\hypertarget{exercise-solution-46}{}
\begin{Shaded}
\begin{Highlighting}[]
\CommentTok{\# 4.}
\NormalTok{survey}\SpecialCharTok{$}\NormalTok{satisfaction\_f }\OtherTok{\textless{}{-}} \FunctionTok{factor}\NormalTok{(survey}\SpecialCharTok{$}\NormalTok{satisfaction, }\AttributeTok{levels =} \DecValTok{1}\SpecialCharTok{:}\DecValTok{5}\NormalTok{,}
                                \AttributeTok{labels =} \FunctionTok{c}\NormalTok{(}\StringTok{"Very dissatisfied"}\NormalTok{, }\StringTok{"Dissatisfied"}\NormalTok{,}
                                           \StringTok{"Neutral"}\NormalTok{, }\StringTok{"Satisfied"}\NormalTok{, }\StringTok{"Very satisfied"}\NormalTok{),}
                                \AttributeTok{ordered =} \ConstantTok{TRUE}\NormalTok{)}
\CommentTok{\# 5.}
\NormalTok{survey}\SpecialCharTok{$}\NormalTok{signup\_date }\OtherTok{\textless{}{-}} \FunctionTok{mdy}\NormalTok{(survey}\SpecialCharTok{$}\NormalTok{signup\_date)}
\NormalTok{survey}\SpecialCharTok{$}\NormalTok{tenure\_days }\OtherTok{\textless{}{-}} \FunctionTok{as.numeric}\NormalTok{(}\FunctionTok{as.Date}\NormalTok{(}\StringTok{"2026{-}11{-}04"}\NormalTok{) }\SpecialCharTok{{-}}\NormalTok{ survey}\SpecialCharTok{$}\NormalTok{signup\_date)}

\CommentTok{\# 6.}
\NormalTok{survey[survey}\SpecialCharTok{$}\NormalTok{renewed, }\FunctionTok{c}\NormalTok{(}\StringTok{"customer"}\NormalTok{, }\StringTok{"annual\_spend"}\NormalTok{)]}
\end{Highlighting}
\end{Shaded}

\begingroup\tiny\color{darkgray}

\begin{verbatim}
  customer annual_spend
1      C01         1200
3      C03           NA
4      C04         2300
6      C06          640
8      C08         1480
\end{verbatim}

\endgroup
\end{frame}

\begin{frame}[fragile]{Bonus: Reference Levels in a Model}
\protect\hypertarget{bonus-reference-levels-in-a-model}{}
\begin{alertblock}{Task}
\protect\hypertarget{task-1}{}
Using the \texttt{Salaries} data:

\begin{enumerate}
\tightlist
\item
  Fit
  \texttt{lm(salary\ \textasciitilde{}\ discipline,\ data\ =\ Salaries)}.
  Which discipline is the reference level, and what does the coefficient
  mean?
\item
  Make the other discipline the reference and refit. What changed, and
  what stayed the same?
\item
  Fit the one-hot version and check that its coefficients equal the
  group means
\end{enumerate}
\end{alertblock}
\end{frame}

\begin{frame}[fragile]{Bonus: Solution}
\protect\hypertarget{bonus-solution}{}
\begin{Shaded}
\begin{Highlighting}[]
\NormalTok{m1 }\OtherTok{\textless{}{-}} \FunctionTok{lm}\NormalTok{(salary }\SpecialCharTok{\textasciitilde{}}\NormalTok{ discipline, }\AttributeTok{data =}\NormalTok{ Salaries)}
\FunctionTok{coef}\NormalTok{(m1)    }\CommentTok{\# intercept = mean of A; coefficient = B minus A}
\end{Highlighting}
\end{Shaded}

\begingroup\tiny\color{darkgray}

\begin{verbatim}
(Intercept) disciplineB 
 108548.431    9480.264 
\end{verbatim}

\endgroup

\begin{Shaded}
\begin{Highlighting}[]
\NormalTok{m2 }\OtherTok{\textless{}{-}} \FunctionTok{lm}\NormalTok{(salary }\SpecialCharTok{\textasciitilde{}} \FunctionTok{fct\_relevel}\NormalTok{(discipline, }\StringTok{"B"}\NormalTok{), }\AttributeTok{data =}\NormalTok{ Salaries)}
\FunctionTok{coef}\NormalTok{(m2)    }\CommentTok{\# intercept = mean of B; coefficient = A minus B}
\end{Highlighting}
\end{Shaded}

\begingroup\tiny\color{darkgray}

\begin{verbatim}
                  (Intercept) fct_relevel(discipline, "B")A 
                   118028.694                     -9480.264 
\end{verbatim}

\endgroup

\begin{Shaded}
\begin{Highlighting}[]
\FunctionTok{all.equal}\NormalTok{(}\FunctionTok{fitted}\NormalTok{(m1), }\FunctionTok{fitted}\NormalTok{(m2))}
\end{Highlighting}
\end{Shaded}

\begingroup\tiny\color{darkgray}

\begin{verbatim}
[1] TRUE
\end{verbatim}

\endgroup

\begin{Shaded}
\begin{Highlighting}[]
\NormalTok{m3 }\OtherTok{\textless{}{-}} \FunctionTok{lm}\NormalTok{(salary }\SpecialCharTok{\textasciitilde{}} \DecValTok{0} \SpecialCharTok{+}\NormalTok{ discipline, }\AttributeTok{data =}\NormalTok{ Salaries)}
\FunctionTok{coef}\NormalTok{(m3)}
\end{Highlighting}
\end{Shaded}

\begingroup\tiny\color{darkgray}

\begin{verbatim}
disciplineA disciplineB 
   108548.4    118028.7 
\end{verbatim}

\endgroup

\begin{Shaded}
\begin{Highlighting}[]
\FunctionTok{tapply}\NormalTok{(Salaries}\SpecialCharTok{$}\NormalTok{salary, Salaries}\SpecialCharTok{$}\NormalTok{discipline, mean)}
\end{Highlighting}
\end{Shaded}

\begingroup\tiny\color{darkgray}

\begin{verbatim}
       A        B 
108548.4 118028.7 
\end{verbatim}

\endgroup
\end{frame}

\hypertarget{sec-summary}{%
\section{Summary}\label{sec-summary}}

\begin{frame}[fragile]{Key Takeaways}
\protect\hypertarget{key-takeaways}{}
\begin{itemize}
\tightlist
\item
  How a variable is stored decides how R treats it: convert numeric
  codes to factors with the data dictionary, and make a deliberate
  choice for Likert items
\item
  Four basic types; factors, Dates, and date-times are built on top of
  them. \texttt{typeof()} shows storage, \texttt{class()} shows behavior
\item
  Use \texttt{is.na()} to find missing values and
  \texttt{na.rm\ =\ TRUE} to ignore them
\item
  The reference level determines what regression coefficients mean: set
  it with \texttt{relevel()}, \texttt{factor(levels\ =\ )}, or
  \texttt{fct\_relevel()}
\item
  Reference and one-hot coding give different coefficients but identical
  predictions; ordered factors use polynomial contrasts
\item
  Vectors, matrices, and arrays hold one type; data frames hold one type
  per column; lists hold anything
\end{itemize}
\end{frame}

\begin{frame}{Next Session and Further Reading}
\protect\hypertarget{next-session-and-further-reading}{}
\textbf{Next session:} operators, functions, and loops

\textbf{Further reading}

\begin{itemize}
\tightlist
\item
  R for Data Science (2e):
  \href{https://r4ds.hadley.nz/factors.html}{Factors}
\item
  R for Data Science (2e):
  \href{https://r4ds.hadley.nz/datetimes.html}{Dates and times}
\item
  Quarto:
  \href{https://quarto.org/docs/presentations/beamer.html}{Beamer
  presentations}
\end{itemize}
\end{frame}



\end{document}
